\documentclass[11pt,oneside,openany]{book}
\usepackage[T1]{fontenc}
\usepackage{lmodern}
\usepackage[letterpaper,margin=1.10in,includeheadfoot]{geometry}
\usepackage{amsmath,amssymb,amsthm,mathtools}
\usepackage{microtype,xcolor,titlesec,fancyhdr,etoolbox,hyperref}
\ifdefined\pdftrailerid
  \pdftrailerid{MathGenome Topics in Homological Algebra 2026-11-20}
\fi
\definecolor{ink}{HTML}{172923}
\definecolor{deepgreen}{HTML}{2B5242}
\definecolor{rulegreen}{HTML}{8A9F92}
\hypersetup{colorlinks=true,linkcolor=deepgreen,urlcolor=deepgreen,citecolor=deepgreen,
  pdftitle={Topics in Homological Algebra},pdfauthor={MathGenome},
  pdfsubject={ALG-10-D01 to ALG-10-D03},bookmarksnumbered=true}
\setcounter{tocdepth}{1}
\setlength{\headheight}{15pt}
\setlength{\emergencystretch}{3em}
\raggedbottom
\titleformat{\chapter}[display]
  {\normalfont\color{ink}}{\color{deepgreen}\large\scshape Chapter \thechapter}{9pt}
  {\huge\bfseries}
\titlespacing*{\chapter}{0pt}{6pt}{24pt}
\titleformat{\section}{\normalfont\Large\bfseries\color{ink}}{\thesection}{0.7em}{}
\pagestyle{fancy}
\fancyhf{}
\fancyhead[L]{\small\sffamily\color{deepgreen}MathGenome}
\fancyhead[R]{\small\nouppercase{\leftmark}}
\fancyfoot[C]{\small\thepage}
\renewcommand{\headrulewidth}{0.35pt}
\renewcommand{\chaptermark}[1]{\markboth{\thechapter.\ #1}{}}
\fancypagestyle{plain}{\fancyhf{}\fancyfoot[C]{\small\thepage}
  \renewcommand{\headrulewidth}{0pt}}
\newtheorem{theorem}{Theorem}[chapter]
\newtheorem{lemma}[theorem]{Lemma}
\newtheorem{proposition}[theorem]{Proposition}
\newtheorem{corollary}[theorem]{Corollary}
\theoremstyle{definition}
\newtheorem{definition}[theorem]{Definition}
\newtheorem{example}[theorem]{Example}
\theoremstyle{remark}
\newtheorem{remark}[theorem]{Remark}
\allowdisplaybreaks[1]

\begin{document}
\frontmatter
\hypersetup{pageanchor=false}
\begin{titlepage}
\thispagestyle{empty}
\thispdfpagelabel{Cover}
\begin{flushleft}
\vspace*{0.68in}
{\large\scshape\color{deepgreen}MathGenome\par}
\vspace{0.35in}
{\color{rulegreen}\rule{0.94\textwidth}{0.65pt}\par}
\vspace{0.42in}
{\color{ink}\fontsize{35}{41}\selectfont\bfseries
Topics in\\[4pt] Homological Algebra\par}
\vspace{0.3in}
{\color{deepgreen}\Large Differential graded categories,\\[4pt]
unbounded derived functors, and descent\par}
\vfill
{\large Mathematical lecture notes\par}
\vspace{0.16in}
{\normalsize ALG-10-D01 through ALG-10-D03\par}
\vspace*{0.45in}
\end{flushleft}
\end{titlepage}
\hypersetup{pageanchor=true}
\begingroup
\small
\makeatletter
\patchcmd{\l@chapter}{\vskip 1.0em \@plus\p@}
  {\vskip 0.45em \@plus 0.25\p@}{}{}
\makeatother
\tableofcontents
\endgroup

\chapter*{Scope and conventions}
\addcontentsline{toc}{chapter}{Scope and conventions}
These notes continue the MathGenome volume \emph{Advanced Homological
Algebra}. They develop differential graded modules and categories, unbounded
derived tensor and internal Hom, and explicit constructions for model
categories, descent and derived affines. The source map records each
numbered result with its hypotheses, location and proof prerequisites.

We use cochain grading and \(H^i(X[s])=H^{i+s}(X)\).
Standard complexes of modules and their tensor products use the Koszul
sign rule. Certain model-category arguments are stated for modules over an
ordinary commutative ring. The derived-affine chapter uses commutative dg
algebras in nonpositive cochain degrees over a characteristic-zero field.
The text identifies the hypotheses whenever it changes base category.
A reference to an earlier MathGenome result by its alphanumeric label
uses the specific proved form of that result.

The theory is developed through finite affine duality, Serre duality
for twists on projective space, explicit faithfully flat descent, and
cotangent and intersection calculations on derived affines. For the
right adjoint to proper pushforward and its perfectness and base-change
hypotheses see \cite[\S\S12--13]{st-duality}. The literature notes
locate homotopy-coherent descent and derived-stack theory.

\mainmatter
\chapter{Differential graded modules and resolutions}
\label{ch:dg-modules}

Fix a commutative ring \(k\). All gradings are cohomological and may be
unbounded in both directions. For a homogeneous element \(x\) we write
\(|x|\) for its degree. The tensor differential on complexes of
\(k\)-modules is
\[
 d(x\otimes y)=dx\otimes y+(-1)^{|x|}x\otimes dy .
\]
The sign rule for this tensor product is also the sign rule for module
actions and for composition. Throughout Chapter~\ref{ch:dg-modules}
we work with right modules. See \cite[\S\S2.2, 3.1--3.2]{keller}
and \cite[\S\S3.3--3.4]{yekutieli} for the background.

\section{Actions, maps, and cones}

\begin{definition}[Differential graded algebras and categories]\label{res:G01}
A differential graded algebra \(A\) is a unital graded \(k\)-algebra
with a differential of degree \(1\) satisfying \(d^2=0\), \(d1=0\), and
\[
 d(ab)=d(a)b+(-1)^{|a|}a\,d(b).
\]
A small dg category \(\mathcal A\) has a set of objects, a complex
\(\mathcal A(a,b)\) for each pair, closed identities of degree \(0\),
and associative unital composition satisfying the same rule with
\(ab\) read as \(a\circ b\). Its category \(H^0(\mathcal A)\) has morphisms
\(H^0(\mathcal A(a,b))\). A dg functor preserves identities,
composition, and the differentials. The one-object construction
identifies dg algebras with one-object dg categories.
\end{definition}
The Leibniz rule makes composition a chain map on the tensor product.
It also shows that composing a boundary with a closed degree-zero
morphism yields a boundary. Hence \(H^0(\mathcal A)\) really is a category.

\begin{definition}[Right modules over a dg category]\label{res:G02}
A right dg \(\mathcal A\)-module \(M\) consists of complexes \(M(a)\)
with associative unital actions
\(M(b)\otimes_k\mathcal A(a,b)\to M(a)\). Write an action as \(x u\).
It satisfies
\[
        d(xu)=d(x)u+(-1)^{|x|}x\,d(u).
\]
The action axioms require \((x u)v=x(u\circ v)\) for composable arrows.
A closed degree-zero map of modules is a family of chain maps
preserving these actions. A quasi-isomorphism of modules is a map
inducing cohomology isomorphisms objectwise.
\end{definition}
The definition for a one-object \(\mathcal A\) is exactly a right dg
module over \(A\). Kernels and cokernels of the degree-zero maps are
objectwise because they are stable under the actions and differentials.

\begin{proposition}[Hom complex and homotopy]\label{res:G03}
There is a dg category of right dg modules whose morphism complex
\(\underline{\operatorname{Hom}}_{\mathcal A}(M,N)\) in degree \(r\)
consists of families of degree-\(r\) graded \(k\)-module maps
\(f_a\) satisfying \(f(xu)=f(x)u\). Its differential is
\[
        \partial f=d_N f-(-1)^r f d_M .
\]
In \(K(\mathcal A)=H^0(\mathrm{dgMod}\text{-}\mathcal A)\), two chain
maps are equal precisely when their difference is \(d_Nh+hd_M\)
for a degree-\(-1\) map \(h\) preserving the right action.
\end{proposition}
\begin{proof}
We use the sign convention with no extra sign in graded
right-linearity. On the expression for \((\partial f)(xu)\),
the two terms involving \(d(u)\) are
\((-1)^{|x|+r}f(x)d(u)\) and
\(-(-1)^{r+|x|}f(x)d(u)\). They cancel, so \(\partial f\) is
again right-linear. Substituting \(\partial f\) into the same
formula gives \(\partial^2f=0\), since \(d_N^2=d_M^2=0\) and
the middle terms have opposite signs. For homogeneous \(g\)
of degree \(s\) we have
\[
 \partial(gf)
 =d_Pgf-(-1)^{s+r}gfd_M
 =(\partial g)f+(-1)^s g(\partial f)
\]
where \(P\) is the target of \(g\). This verifies the dg category
structure. The statement about \(K(\mathcal A)\) follows by taking
the degree-zero cycles and boundaries in this very complex.
\end{proof}

\begin{proposition}[Shift, cone, and the acyclicity test]\label{res:G04}
For \(s\in\mathbb Z\), put \(M[s]^i=M^{i+s}\) with differential
\((-1)^s d_M\) and the same right action. For a closed degree-zero
map \(f:M\to N\), its cone is \(N\oplus M[1]\) with differential
\[
 d(y,x)=(d_N y+f(x),-d_M x).
\]
The map \(f\) is a quasi-isomorphism if and only if its cone is
acyclic. The cone and shift act objectwise. The homotopy category
\(K(\mathcal A)\) is triangulated with the triangles isomorphic to
these cone triangles.
\end{proposition}
\begin{proof}
After a shift the Leibniz term acquires the sign
\((-1)^{s+|x|_M}=(-1)^{|x|_{M[s]}}\). The shift thus respects the
right action. On a homogeneous map of degree \(r\), multiplication
by \((-1)^{sr}\) upon reindexing gives the dg shift functor, as
substitution into \(\partial\) verifies. The cone differential
squares to zero since \(f\) is a chain map, and the term \(f(xu)\)
equals \(f(x)u\). Thus the cone is a dg module. The short exact
sequence \(0\to N(a)\to C(f)(a)\to M(a)[1]\to 0\) identifies
its connecting homomorphism with \(H(f(a))\) by lifting a cycle
as \((0,x)\). Exactness gives the equivalence.

The cone proof of the triangulated axioms in
\cite[K01--K06]{ha-core} applies with right-linear components.
More precisely, the comparison of homotopic maps, the rotation,
the map of cones obtained from a specified homotopy of squares,
and the octahedral maps there use only identities, compositions,
sums, shifts and the cone matrix. All these formulas preserve
right-linearity by the verifications above. The component
homotopies in the cone proof are also right-linear, since they
are matrices formed from the same operations on the input
homotopies. Therefore these are triangles and the same comparison
maps inside \(K(\mathcal A)\).
\end{proof}

\section{Representables and free cells}

\begin{lemma}[Enriched Yoneda for right modules]\label{res:G05}
For \(a\in\mathcal A\) put \(h_a(b)=\mathcal A(b,a)\) with the
right action by composition. Evaluation at \(1_a\) is a natural
isomorphism of complexes
\[
     \underline{\operatorname{Hom}}_{\mathcal A}(h_a,M)
        \longrightarrow M(a).
\]
The assignment \(a\mapsto h_a\) is a fully faithful dg functor.
\end{lemma}
\begin{proof}
For \(m\in M(a)^r\), its inverse map has components
\(u\mapsto m u\) on \(\mathcal A(b,a)\). This is graded right-linear.
Any graded map out of \(h_a\) must have that formula since
\(u=1_a u\). The terms involving \(d u\) in its Hom differential
cancel by the Leibniz rule, leaving the map \(u\mapsto (d m)u\).
On representables the Yoneda map takes composition to composition
as \((m u)v=m(u\circ v)\). These checks give the asserted dg
functor and full faithfulness.
\end{proof}

A free graded cell of degree \(n\) at \(a\) is a copy of
\(h_a[-n]\) with a named element \(e\) corresponding to \(1_a\).
One can adjoin it to a dg module \(P\) by setting \(de=z\) for
a chosen cycle \(z\in P(a)^{n+1}\). On products the
differential must be
\[
                 d(eu)=z u+(-1)^n e\,d u.
\]
This formula squares to zero. Indeed \(d(z u)=(-1)^{n+1}z\,du\),
which cancels the contribution from differentiating the other
term. The underlying graded module is the direct sum with the
new cell. Its quotient by \(P\) is \(h_a[-n]\). A semi-free module
here has an exhaustive filtration
\(0=P_{-1}\subset P_0\subset P_1\subset\cdots\)
obtained by adjoining arbitrary sets of such cells, each with
differential in the preceding stage. All sums are direct sums of
graded modules, so each element involves finitely many cells.

\begin{theorem}[Unbounded semi-free resolution]\label{res:G06}
Every dg \(\mathcal A\)-module \(M\) admits a degreewise surjective
quasi-isomorphism \(P\to M\) from a semi-free module.
\end{theorem}
\begin{proof}
Put in stage zero a cell of degree \(n\) with zero differential
for every cycle in \(M(a)^n\) at every object and every degree.
Map the generator to that cycle. Suppose \(p_j:P_j\to M\) is
defined. For each \(z\in P_j(a)^n\) with \(dz=0\), and each
\(x\in M(a)^{n-1}\) satisfying \(dx=p_j(z)\), adjoin in stage \(j+1\)
a generator \(e\) of degree \(n-1\) with \(de=z\) and send it to \(x\).
The displayed cell formula and the equation on the proposed image
make this a chain map. The indexing uses only sets because
\(\mathcal A\) is small.

Any class in \(H^n(M(a))\) is represented by a cycle from stage zero.
A cycle of \(P(a)\) whose image is a boundary is already a cycle
at a finite stage because the filtrations consist of injections
and each element has finite support. At the next stage a cell
with a chosen primitive of its image kills that cycle. Hence
\(H^n(P(a))\to H^n(M(a))\) is both surjective and injective.
Finally, for an arbitrary \(x\in M(a)^n\) the cycle \(dx\) was
represented in stage zero. The new cell attached for this cycle
and \(x\) gives a degree-\(n\) preimage at stage one. This proves
degreewise surjectivity.
\end{proof}

\chapter{Derived dg categories and Morita theory}
\label{ch:dg-categories}

The representable resolutions above let us compute the derived
category without any boundedness condition. Its compact objects
recover the finite shifts, cones, and retracts of the representables.
We then use the signed tensor--Hom adjunction to compare derived
categories of different dg categories.

\section{Homotopy-projectivity and the derived category}

\begin{theorem}[Semi-free modules are \(K\)-projective]\label{res:G07}
A dg module \(P\) is called \(K\)-projective if
\(\underline{\operatorname{Hom}}_{\mathcal A}(P,N)\) is acyclic for
every acyclic \(N\). Every semi-free module is \(K\)-projective.
The class of \(K\)-projective modules is closed under shifts,
cones of closed degree-zero maps, and arbitrary direct sums.
\end{theorem}
\begin{proof}
Let \(N\) be acyclic and let \(f:P\to N\) be a graded map of
degree \(r\) with \(\partial f=0\), where \(P\) is semi-free.
Construct a graded map \(h\) of degree \(r-1\) on successive
cell stages with \(\partial h=f\). At a new generator \(e\)
of degree \(n\) at \(a\), take \(z=de\) in the previous stage.
The requirement is
\[
          d_N h(e)= f(e)-(-1)^r h(z).
\]
The right side is closed because
\(d_N f(e)=(-1)^r f(z)\), and the already constructed homotopy
gives \(d_N h(z)=f(z)\) since \(dz=0\). Acyclicity of \(N(a)\)
gives a primitive, which we choose as \(h(e)\). Extending
right-linearly over the new cells is well-defined, and verification
on generators suffices by Result~\ref{res:G03}. Taking the union yields
\(\partial h=f\) at every degree. This proves \(K\)-projectivity.

Shifts preserve the condition by reindexing Hom complexes with
the shift signs of Result~\ref{res:G04}. For a cone \(C\) of
\(P_1\to P_2\), the degreewise split sequence
\(P_2\to C\to P_1[1]\) gives a short exact sequence of Hom complexes
into \(N\), in the reverse order, since its splitting exists
in graded \(\mathcal A\)-modules. The end complexes are acyclic.
For a sum of \(P_i\), Hom becomes the product of their Hom
complexes. A product of acyclic complexes of modules is acyclic
by choosing a primitive in every component of a cycle.
\end{proof}

\begin{theorem}[Computing the derived category by resolutions]\label{res:G08}
Let \(D(\mathcal A)\) be the localization of \(K(\mathcal A)\)
at the objectwise quasi-isomorphisms. The full homotopy subcategory
of \(K\)-projective modules maps equivalently onto \(D(\mathcal A)\).
If \(P\to M\) is a \(K\)-projective replacement, then
\[
 \operatorname{Hom}_{D(\mathcal A)}(M,N)
      \cong H^0\underline{\operatorname{Hom}}_{\mathcal A}(P,N).
\]
Cones of maps of \(K\)-projective modules represent distinguished
triangles, and coproducts of arbitrary derived objects are
represented by direct sums of their modules.
\end{theorem}
\begin{proof}
For a quasi-isomorphism \(u:N'\to N\) and \(K\)-projective \(P\),
\(\underline{\operatorname{Hom}}_{\mathcal A}(P,u)\) is a
quasi-isomorphism. Its cone is identified with Hom into the cone
of \(u\) using the Hom and shift formulas, hence is acyclic.
An acyclic \(K\)-projective module is contractible since its
identity represents a class in \(H^0\) of its own acyclic Hom
complex. By the closure in Result~\ref{res:G07}, the cone of a
quasi-isomorphism between \(K\)-projective modules is contractible,
so that quasi-isomorphism is a homotopy equivalence.

For every \(M\) choose \(p_M:P_M\to M\) as in Result~\ref{res:G06},
taking the identity when \(M\) is already \(K\)-projective.
For a homotopy class \(M\to N\), the isomorphism
\(\operatorname{Hom}_{K(\mathcal A)}(P_M,P_N)\to
  \operatorname{Hom}_{K(\mathcal A)}(P_M,N)\)
induced by \(p_N\) gives a unique lifted homotopy class
\(P_M\to P_N\). Uniqueness proves the composition law and
naturality of the \(p_M\). If the original map is a quasi-isomorphism,
so is its lift by the comparison square, and this lift is
invertible in the homotopy category. The lift functor therefore
factors through the localization. The comparisons \(p_M\) give a
natural isomorphism from its composite with inclusion to the
identity on the localization. In the other order the composite
is the identity on the \(K\)-projective subcategory. This proves
the equivalence and the stated Hom formula, since \(N\) can be
compared with \(P_N\) after applying Hom from \(P_M\).

The \(K\)-projective homotopy subcategory is triangulated by
Results~\ref{res:G04} and~\ref{res:G07}. It gives the cone
triangulation on the localization. Indeed a closed map between
arbitrary modules is represented by the lifted map between
replacements. A comparison with its original cone exists from
the degree-\(-1\) homotopy in the comparison square and the
usual cone matrix in \cite[K05]{ha-core}. It is a quasi-isomorphism
by the long exact sequences for the two cones. Direct sums of
\(K\)-projective modules are \(K\)-projective, and direct sums preserve
objectwise quasi-isomorphisms. Hom from a sum of replacements is
the product of the corresponding Hom groups, as \(H^0\) commutes
with products of complexes of modules. Hence the sum of the
chosen replacements computes the coproduct, and its map to the
sum of original modules is a quasi-isomorphism.
\end{proof}

\section{Compactness and the pretriangulated hull}

\begin{lemma}[The telescope of a sequence of strict maps]\label{res:G09}
Let \(M_0\to M_1\to\cdots\) be a sequence of dg modules and closed
degree-zero maps. On \(U=\bigoplus_{j\geq0}M_j\) let \(s\) send
the \(j\)-th summand to the \((j+1)\)-st by the transition.
Then \(\operatorname{Cone}(1-s)\to\varinjlim_j M_j\) is an
objectwise quasi-isomorphism. If \(C\) is compact in
\(D(\mathcal A)\), any map from \(C\) to this telescope in
\(D(\mathcal A)\) comes from a single \(M_j\). In fact its group
of maps is the colimit of the groups of maps into \(M_j\).
\end{lemma}
\begin{proof}
The map \(1-s\) is injective in each degree by considering the
first nonzero component of a finite-support element. Its cokernel
imposes exactly the relations identifying an element and its
successor, which presents the colimit. For any injective dg map
\(U'\to U''\), the natural map from its cone to its cokernel
has kernel the cone of the identity of \(U'\). That cone is
contractible by the right-linear homotopy
\((y,x)\mapsto(0,y)\). This proves the claim for \(1-s\).

Compactness means that \(\operatorname{Hom}(C,-)\) takes arbitrary
sums to direct sums, and this also applies to the objects \(M_j[1]\).
In the Hom exact sequence of the cone of \(1-s\), the map
\(1-s\) on the direct sum of the groups
\(\operatorname{Hom}(C,M_j[1])\) is injective by the same first
component argument. Thus \(\operatorname{Hom}(C,\operatorname{Cone}(1-s))\)
is the cokernel of \(1-s\) on the unshifted Hom groups. This is
their colimit. An element in a sequential colimit of groups
can be represented at a single stage.
\end{proof}

\begin{theorem}[Representables are compact generators]\label{res:G10}
Each \(h_a\) is a compact object of \(D(\mathcal A)\). For every
\(N\) and \(n\in\mathbb Z\),
\[
       \operatorname{Hom}_{D(\mathcal A)}(h_a[-n],N)
           \cong H^n(N(a)).
\]
The smallest full triangulated subcategory closed under
isomorphisms and coproducts and containing the \(h_a\) is all of
\(D(\mathcal A)\). In particular the shifted representables detect
the zero object.
\end{theorem}
\begin{proof}
The Hom formula follows from enriched Yoneda and
Results~\ref{res:G05} and~\ref{res:G08}. Explicitly a map out of
the shifted cell is determined by a degree-\(n\) element, closed
for a chain map, and homotopy identifies classes modulo boundaries.
Since cohomology of modules commutes with direct sums, the formula
also proves compactness by the coproduct description of
Result~\ref{res:G08}.

For generation resolve \(N\) by a semi-free module \(P\). Each
\(P_j/P_{j-1}\) is a direct sum of shifted representables, and
\(P_{j-1}\to P_j\to P_j/P_{j-1}\) is part of a triangle by the
cone-to-cokernel argument of Result~\ref{res:G09}. Therefore a
subcategory with the stated closure contains each \(P_j\).
It contains the telescope, which represents \(P\) by
Result~\ref{res:G09}, and hence contains \(N\). Alternatively,
vanishing of the displayed Hom groups says that \(N\) is
objectwise acyclic, so \(N\) is zero in the localization.
\end{proof}

\begin{theorem}[Compact objects and finite cells]\label{res:G11}
An object of \(D(\mathcal A)\) is compact if and only if it is
a retract in \(D(\mathcal A)\) of a semi-free module with only
finitely many cells.
\end{theorem}
\begin{proof}
Finite-cell objects are obtained by finite sums, shifts, and
cones from representables. Compact objects are closed under
those operations by the exact Hom sequences of triangles and
exactness of direct sums of groups. A retract is compact because
the coproduct comparison map for a retract is a retract of the
same map for the original object.

Conversely let \(C\) be compact and choose a semi-free
representative \(P\) of \(C\). By Result~\ref{res:G09} its identity
factors in \(D(\mathcal A)\) through one finite stage \(P_j\) of
the cell filtration, which may itself contain infinitely many
cells. At a fixed stage \(j\), call a finite subset of cells
closed if it contains all the cells occurring in the differential of
every cell in it. Every finite subset can be enlarged to a finite
closed subset by recursively adding the finite differential
supports, which lie at strictly earlier stages. Finite closed
subsets form a filtered system under unions. For such a subset
\(E\), let \(P_{j,E}\) be its finite-cell submodule.

For every \(t\) there is an isomorphism
\[
 \varinjlim_E\operatorname{Hom}_{D(\mathcal A)}(C,P_{j,E}[t])
     \ \cong\ \operatorname{Hom}_{D(\mathcal A)}(C,P_j[t]).
\]
Prove it by induction on \(j\). For stage zero it is compactness
on a sum of shifted representables, since each element of the sum
has finite support. At stage \(j\), the first \(j-1\) stages for a
finite \(E\) form \(P_{j-1,E\cap P_{j-1}}\), and the quotient is a
finite sum of cells from stage \(j\). The associated triangles,
including all shifts, give a comparison of long exact Hom sequences
after taking the filtered colimit. Filtered colimits of groups
are exact, since a finite relation between representatives holds
at a sufficiently late common stage. The group maps for the
lower-stage terms are isomorphisms by induction. The system of
intersections is cofinal among finite closed subsets of the
lower stage, since any such subset can itself be used as \(E\).
The quotient term maps are isomorphisms by compactness, since
each finite subset of top-stage cells can be enlarged to a finite
closed subset. Exactness on the two sides implies the middle
group maps are isomorphisms by the five-term diagram argument.
This proves the displayed comparison.

Consequently the factorization of the identity through \(P_j\)
factors further through a finite-cell \(P_{j,E}\). Its composite
back to \(C\) is the identity, as required.
\end{proof}

\begin{proposition}[The pretriangulated hull and perfect objects]\label{res:G12}
There is a small full dg category \(\operatorname{pretr}(\mathcal A)\)
of finite-cell semi-free modules that contains the Yoneda image,
represents all finite-cell semi-free modules up to dg isomorphism,
and is closed under zero, shifts, finite sums, and cones of closed
degree-zero maps. Its homotopy category embeds as a triangulated
subcategory of \(D(\mathcal A)\). The perfect subcategory,
obtained by also taking retracts in \(D(\mathcal A)\), is exactly
the subcategory of compact objects. The Yoneda map
\(\mathcal A\to\operatorname{pretr}(\mathcal A)\) is a quasi-equivalence
precisely when all finite-cell objects are represented, up to
homotopy, by objects of \(\mathcal A\).
\end{proposition}
\begin{proof}
The cone of a map from \(P\) to \(Q\) has graded module
\(Q\oplus P[1]\). Order all cells of \(Q\) first and then those
of \(P[1]\) in their existing order of stages. The shifted
\(P\)-differential lowers these latter stages, and the cone's
off-diagonal term lies in \(Q\). This orders the cone as a
finite-cell module. Shifts preserve the decreasing property of
the cell differential, and finite sums admit an ordering by
the original stages. Conversely attaching one cell of degree
\(n\) with boundary \(z\) is the cone of
\(h_a[-n-1]\to P\) carrying the generator to \(z\), by the cone
formula. Finite-cell modules are therefore exactly the finite
constructions generated by the representables in this model,
up to choices of graded cell presentations. For smallness,
start with the set of written representables and zero. At
each finite stage adjoin the written shifts, finite sums and
cones of every closed degree-zero map among the present objects,
using the formulas of Result~\ref{res:G04}. These data range
over sets because each Hom is a set. The union over stages
is a set closed under the stated operations. Take the full
dg subcategory on this set.

All finite-cell modules are \(K\)-projective. By
Result~\ref{res:G08} their homotopy morphisms compute the
derived morphisms, and by the cone calculation the full image
is triangulated. Result~\ref{res:G11} gives the assertion about
perfect objects. A dg functor is a quasi-equivalence if it
induces quasi-isomorphisms on all morphism complexes and is
essentially surjective on \(H^0\). Yoneda is already an
isomorphism on the morphism complexes, so the condition is
exactly the stated representability up to homotopy. In this
case the triangulation transports to \(H^0(\mathcal A)\).
\end{proof}

\section{Derived tensor functors and Morita equivalence}

Let \(\mathcal A,\mathcal B\) be small dg categories over \(k\).
An \(\mathcal A\)-\(\mathcal B\) dg bimodule will be described by a
dg functor \(T:\mathcal A\to \mathrm{dgMod}\text{-}\mathcal B\).
In particular \(u:a\to b\) acts as \(T(u):T(a)\to T(b)\) of the
same degree as \(u\), in the Hom convention of Result~\ref{res:G03}.

\begin{proposition}[The signed bimodule adjunction]\label{res:G13}
For a right dg \(\mathcal A\)-module \(M\), \(M\otimes_\mathcal A T\)
is the quotient of \(\bigoplus_a M(a)\otimes_k T(a)\) by the
homogeneous balancing relations
\((m u)\otimes t=m\otimes T(u)t\).
For a right dg \(\mathcal B\)-module \(N\), put
\(G(N)(a)=\underline{\operatorname{Hom}}_\mathcal B(T(a),N)\)
with right \(\mathcal A\)-action by precomposition. There is a
natural isomorphism of complexes
\[
 \underline{\operatorname{Hom}}_\mathcal B(M\otimes_\mathcal A T,N)
    \cong \underline{\operatorname{Hom}}_\mathcal A(M,G(N)).
\]
If each \(T(a)\) is \(K\)-projective, \(G\) sends quasi-isomorphisms
to quasi-isomorphisms. Computing the tensor on a \(K\)-projective
replacement yields an exact left adjoint
\(L_T:D(\mathcal A)\to D(\mathcal B)\) to \(G\) on derived categories.
Moreover \(L_T(h_a)\cong T(a)\).
\end{proposition}
\begin{proof}
The differential on the direct sum is the tensor differential.
It preserves the balancing relation because its derivative splits
into the three relations for \(dm,u,t\), \(m,du,t\), and \(m,u,dt\),
with coefficients \(1,(-1)^{|m|}\), and \((-1)^{|m|+|u|}\).
This uses
\(d(T(u)t)=T(du)t+(-1)^{|u|}T(u)dt\), which says that \(T\) is dg.
The right \(\mathcal B\)-action on the second tensor factor
respects balancing, since each \(T(u)\) is graded right-linear.
The tensor differential obeys the right Leibniz rule for this
action. For the right \(\mathcal A\)-module \(G(N)\), write
\(g u=g\circ T(u)\). The Hom composition rule of
Result~\ref{res:G03} gives its Leibniz rule.

If \(F:M\otimes_\mathcal A T\to N\) has degree \(r\), send it to
the graded map \(m\mapsto (t\mapsto F(m\otimes t))\). Balancing
is precisely \(\mathcal A\)-linearity of the curried map, and its
values are \(\mathcal B\)-linear maps. Conversely these two
linearity conditions define the balanced uncurried map.
For the differential, evaluating the curried differential
on \(m,t\) gives
\[
 d_N F(m\otimes t)-(-1)^r F(dm\otimes t)
       -(-1)^{r+|m|}F(m\otimes dt),
\]
which is the currying of the uncurried Hom differential.
This proves the isomorphism and its compatibility with closed
maps. The unit and counit are \(m\mapsto[t\mapsto m\otimes t]\)
and \(g\otimes t\mapsto g(t)\). Their composites on evaluated
tensors and maps are identities.

Now suppose the \(T(a)\) are \(K\)-projective. The formula of
Result~\ref{res:G08} for applying Hom from such modules shows
that \(G\) takes acyclics to acyclics and quasi-isomorphisms to
quasi-isomorphisms. For \(K\)-projective \(P\), the adjunction
of complexes shows \(P\otimes_\mathcal A T\) is \(K\)-projective
because its Hom into an acyclic \(N\) is Hom from \(P\) into the
acyclic \(G(N)\). The tensor takes homotopy equivalences to
homotopy equivalences by tensoring the homotopy with the identity.
Thus compute \(L_T\) on \(K\)-projective representatives.
The comparison of derived Hom groups in Result~\ref{res:G08}
now gives the derived adjunction. Tensoring commutes with shifts
and cones by the differential formulas, so \(L_T\) is exact.
For \(G\), the cone and shift in the target of Hom are likewise
the cone and shift of Hom, hence \(G\) is exact on the derived
categories. Finally \(h_a\otimes_\mathcal A T\to T(a)\) is
given by \(u\otimes t\mapsto T(u)t\), with inverse
\(t\mapsto 1_a\otimes t\). Balancing proves both composites are
identities and the maps commute with differentials.
\end{proof}

\begin{theorem}[A derived Morita criterion]\label{res:G14}
Assume each \(T(a)\) is \(K\)-projective, as in
Result~\ref{res:G13}. The derived adjunction is an equivalence
if the following conditions hold.
\begin{enumerate}
\item Each \(T(a)\) is compact in \(D(\mathcal B)\).
\item For every \(a,b\), the map of complexes
\[
  \mathcal A(a,b)\longrightarrow
     \underline{\operatorname{Hom}}_\mathcal B(T(a),T(b))
\]
induced by \(T\) is a quasi-isomorphism.
\item The shifted objects \(T(a)[n]\) detect zero in \(D(\mathcal B)\).
\end{enumerate}
Conversely, if \(L_T\) is an equivalence, these conditions hold.
\end{theorem}
\begin{proof}
Compactness and Result~\ref{res:G08} show that the natural map
\(\bigoplus_i G(N_i)\to G(\bigoplus_i N_i)\) is a quasi-isomorphism.
Indeed its cohomology at each \(a,n\) is the coproduct comparison
for maps from \(T(a)[-n]\), which is compact since \(T(a)\) is.
The functor \(L_T\) commutes with coproducts by the adjunction
or the tensor formula. Both functors are exact by
Result~\ref{res:G13}. In particular the full subcategory where
the unit \(X\to G(L_T(X))\) is an isomorphism is closed under
shifts, cones, isomorphisms, and coproducts. Closure under cones
follows by applying the Hom exact sequences to the comparison
of triangles. The unit on \(h_b\), evaluated at \(a\), is
exactly the map in condition (2), by the formula for the unit and
the evaluation \(h_b\otimes_\mathcal A T\cong T(b)\).
Thus this subcategory is all of \(D(\mathcal A)\) by
Result~\ref{res:G10}.

For \(Y\in D(\mathcal B)\) apply \(G\) to the counit
\(L_T G(Y)\to Y\). The triangular identity identifies this as
the inverse of the unit on \(G(Y)\). Condition (3) implies \(G\)
detects zero by computing its cohomology at each \(a\) using
maps from the shifts of \(T(a)\). Since it is exact it detects
isomorphisms by applying the zero detection to the cone.
Thus the counit is an isomorphism. This proves equivalence.

Conversely \(L_T\) carries each \(h_a\) to \(T(a)\).
An exact equivalence preserves coproducts and the compact
generation and detection properties in Result~\ref{res:G10}.
Full faithfulness on all shifted representables identifies the
cohomology in every degree of the map in condition (2) with
an isomorphism of derived morphism groups. This identification
uses the Yoneda evaluation on the source and the
\(K\)-projectivity of \(T(a)\) on the target. Hence that map is
a quasi-isomorphism.
\end{proof}

\begin{corollary}[Invariance under quasi-equivalence]\label{res:G15}
A quasi-equivalence \(F:\mathcal A\to\mathcal B\) induces an
equivalence of derived categories by extension of scalars, with
restriction along \(F\) as inverse.
\end{corollary}
\begin{proof}
Use \(T(a)=h_{F(a)}\) with its Yoneda functoriality. Enriched Yoneda
identifies \(G\) with restriction. Each \(T(a)\) is a compact
\(K\)-projective module. The quasi-isomorphism condition on
morphism complexes of \(F\) is condition (2) of
Result~\ref{res:G14}. Every \(b\) is isomorphic to some \(F(a)\)
in \(H^0(\mathcal B)\), so every \(h_b\) is isomorphic to
some \(T(a)\) in the derived category. The detection assertion
of Result~\ref{res:G10} gives condition (3).
\end{proof}

\begin{corollary}[An endomorphism dg algebra of a generator]\label{res:G16}
Let \(P\) be a compact \(K\)-projective right \(\mathcal B\)-module
whose shifts detect zero. Then \(D(\mathcal B)\) is equivalent
to \(D(\underline{\operatorname{End}}_\mathcal B(P))\).
Also the inclusion of \(\mathcal A\) into its pretriangulated
hull induces an equivalence of derived categories.
\end{corollary}
\begin{proof}
Regard \(\underline{\operatorname{End}}_\mathcal B(P)\) as a
one-object dg category acting on \(P\). Its map in condition (2)
of Result~\ref{res:G14} is the identity. The remaining conditions
are the hypotheses on \(P\), proving the first assertion.
For the hull, use Yoneda in \(\operatorname{pretr}(\mathcal A)\)
on each object \(h_a\) in the image of \(\mathcal A\).
The induced maps of morphism complexes are isomorphisms.
A finite-cell object of the hull is formed by shifts and cones
from the \(h_a\). The representable module of that finite-cell
object belongs to the triangulated subcategory generated by
the representables of \(h_a\). In fact Hom into a shift is a
shift of Hom, and Hom into a cone is the corresponding cone
objectwise, by the differential of Result~\ref{res:G03}.
Thus the representables of \(h_a\) detect zero in
\(D(\operatorname{pretr}(\mathcal A))\) by Result~\ref{res:G10}
and exact Hom from triangles. Apply Result~\ref{res:G14} again.
\end{proof}

For a concrete application take a field \(k\) and \(P=k^m\)
over \(k\) with \(m\geq2\). It is \(K\)-projective, compact,
and its shifts detect zero since \(k\) is a direct summand.
Its endomorphism dg algebra is \(\operatorname{Mat}_m(k)\)
in degree zero. Thus \(D(k)\) and \(D(\operatorname{Mat}_m(k))\)
are equivalent by Result~\ref{res:G16}. Any quasi-equivalence
of the corresponding one-object dg categories would give an
isomorphism of their degree-zero endomorphism \(k\)-algebras,
which have different dimensions. Even a zigzag of
quasi-equivalences induces equivalences of the \(H^0\) categories
and would give the same impossible isomorphism. This illustrates
the distinction between quasi-equivalence and derived Morita
equivalence. The criterion above concerns an explicit bimodule
functor, so its conclusion also records the dg structure used
to induce the equivalence.

The dg-category localization at Morita equivalences and the
comparison of arbitrary enhancements involve further structure.
See \cite[\S\S3.8, 4.5]{keller} for these developments.

\chapter{Unbounded derived tensor and internal Hom}
\label{ch:unbounded}

Throughout the chapter \(R\) is a commutative ring. Complexes are
cochain complexes, with \(X[s]^i=X^{i+s}\) and
\(d_{X[s]}=(-1)^s d_X\). Modules are placed in degree zero. We use the
construction of bounded-below injective replacements and the homotopy
calculus of \cite[R01, R03, R04]{ha-core}. In particular, no assumption
of Noetherianity is needed in this chapter.

\section{Complexes and free cells}

\begin{definition}[Tensor, internal Hom, and good models]\label{res:U01}
The tensor complex is totalized by direct sums, with
\(d(x\otimes y)=d x\otimes y+(-1)^{|x|}x\otimes d y\).
The internal Hom complex is
\[
 \underline{\mathrm{Hom}}_R(X,Y)^q
   =\prod_i\mathrm{Hom}_R(X^i,Y^{i+q}),\qquad
 \partial f=d_Y f-(-1)^q f d_X .
\]
A complex \(P\) is K-projective when
\(\underline{\mathrm{Hom}}_R(P,A)\) is acyclic for every acyclic \(A\).
A complex \(F\) is K-flat when \(A\otimes_R F\) is acyclic for every
acyclic \(A\). A complex \(I\) is K-injective when
\(\underline{\mathrm{Hom}}_R(A,I)\) is acyclic for every acyclic \(A\).
\end{definition}

Expanding the two differentials proves that their squares are zero.
For the Hom differential, the two terms containing \(d_Y f d_X\)
cancel. Cycles in degree zero are chain maps, and degree-zero
boundaries are null-homotopic maps. The vanishing condition in all
degrees is thus the usual orthogonality to all shifts of acyclic
complexes. Direct sums in the tensor and products in the Hom are
essential when both arguments are unbounded. For example the
degree-\(q\) tensor--Hom adjunction is the correspondence
\[
 u(x\otimes y)\ \longleftrightarrow\ [x\longmapsto(y\longmapsto
 u(x\otimes y))].
\]
For \(u\) of degree \(q\), each inner map has degree \(q+|x|\).
Expanding \(\partial u\) gives the three terms acting on the target,
on \(x\), and on \(y\), with coefficients \(1,-(-1)^q\),
and \(-(-1)^{q+|x|}\), respectively. This is exactly the Hom
differential on the curried expression. The constructions and the
adjunction carry chain homotopies to chain homotopies.

\begin{theorem}[A semifree replacement without bounds]\label{res:U02}
Every complex \(X\) admits a degreewise surjective quasi-isomorphism
\(F\to X\) such that
\(F\) has an increasing filtration by subcomplexes
\(0\subset F_0\subset F_1\subset\cdots\) with union \(F\).
Each inclusion is split on the underlying graded modules and the
successive quotients are sums of shifts of \(R\).
The complex \(F\) is both K-projective and K-flat.
\end{theorem}
\begin{proof}
View \(R\) as a one-object dg category in degree zero.
Its representable module is the complex \(R[0]\).
Theorems~\ref{res:G06} and \ref{res:G07} give the
graded-split filtration, quasi-isomorphism, and K-projectivity.

Tensor the filtration with \(A\). Splitting on graded modules makes
each successive sequence exact as complexes. Each quotient is a
direct sum of shifts of \(A\), hence acyclic. The finite-stage tensor
complexes are therefore acyclic by induction and their filtered
union remains so, since filtered colimits are exact for modules.
Thus \(F\) is K-flat.
\end{proof}

\begin{proposition}[Invariance and bounded tests]\label{res:U03}
The functors \(P\otimes_R-\), \(\underline{\mathrm{Hom}}_R(P,-)\),
and \(\underline{\mathrm{Hom}}_R(-,I)\) preserve quasi-isomorphisms
when \(P\) is K-flat, \(P\) is K-projective, and \(I\) is
K-injective, respectively. A bounded-above complex of flat modules
is K-flat. Every K-projective complex is K-flat.
\end{proposition}
\begin{proof}
The cone of a quasi-isomorphism is acyclic. Tensor takes a cone to
the cone of the tensor map, up to the canonical shift sign.
In the component involving \(x\in P^i\) and the shifted input the
comparison is \((-1)^i\), as is seen by expanding the differential
in Definition~\ref{res:U01}. Internal Hom in either variable
gives the corresponding cone up to shift, since the cone uses
a graded finite direct sum and its differential contains the
original map as off-diagonal entry. Thus the three acyclicity
conditions imply the asserted invariance.

For a bounded-above complex of flats, first cut off the degrees
below \(m\) by the subcomplex consisting of its terms in degrees
at least \(m\). This is a bounded complex filtered by degree tails,
whose one-term quotients are shifts of flat modules. Tensoring with
an acyclic complex preserves those graded-split exact sequences,
and each one-term quotient gives an acyclic tensor by flatness.
As \(m\) decreases the union is the original complex, so exactness
of filtered colimits proves K-flatness. Finally, apply
Theorem~\ref{res:U02} to a K-projective complex. The quasi-isomorphism
from its semifree replacement is between K-projectives, hence is a
homotopy equivalence by \cite[R04]{ha-core}. Tensor with an acyclic
complex preserves that homotopy equivalence. 
\end{proof}

\section{Injective replacements from a tower}

\begin{lemma}[Strict extension to an injective model]\label{res:U04}
Let \(i:X\hookrightarrow J\) be a termwise monic quasi-isomorphism.
If \(I\) is a K-injective complex with injective terms, every
chain map \(f:X\to I\) extends to a chain map \(J\to I\).
\end{lemma}
\begin{proof}
Extend the components first, using injectivity, to obtain a graded
map \(u:J\to I\) of degree zero. Since \(u i=f\), the degree-one
map \(\partial u\) vanishes on \(X\) and factors through \(J/X\)
as a cycle in \(\underline{\mathrm{Hom}}_R(J/X,I)\).
The quotient \(J/X\) is acyclic by the long exact sequence.
K-injectivity makes this cycle the differential of a degree-zero
map \(v:J/X\to I\). Subtract \(v\) composed with the quotient
map from \(u\). The result extends \(f\) and is closed.
\end{proof}

\begin{lemma}[Countable towers of complexes]\label{res:U05}
For a tower of complexes \(A_0\leftarrow A_1\leftarrow\cdots\)
whose transition maps are degreewise surjective, there is a short
exact sequence of complexes
\[
 0\longrightarrow\varprojlim_n A_n\longrightarrow
 \prod_{n\geq0} A_n
 \xrightarrow{\ 1-s\ }\prod_{n\geq0} A_n
 \longrightarrow0,
\qquad
 s((a_n)_n)=(t_n(a_{n+1}))_n .
\]
A degreewise surjective tower of acyclic complexes has acyclic
limit. A map between degreewise surjective towers that is a
quasi-isomorphism at each stage gives a quasi-isomorphism of
limits. If all \(A_n\) are K-injective and transitions split
in each degree, then the limit is K-injective.
\end{lemma}
\begin{proof}
In each degree the kernel is the compatible families.
To solve \(a_n-t_n(a_{n+1})=b_n\), choose \(a_0\) and then
lift \(a_n-b_n\) recursively. Products of modules are exact,
since surjectivity can be tested in each component. They
therefore commute with taking cohomology of complexes.
The first two consequences follow from the long cohomology
sequences of the displayed short exact sequence.
For the last one, if \(C\) is acyclic pass to the tower
\(\underline{\mathrm{Hom}}_R(C,A_n)\). Its terms are acyclic,
its transition maps are surjective by the degreewise splittings,
and its limit is \(\underline{\mathrm{Hom}}_R(C,\varprojlim A_n)\)
since limits commute with products and with Hom in the target.
Apply the first consequence.
\end{proof}

\begin{theorem}[Unbounded K-injective replacement over a ring]\label{res:U06}
Every complex \(X\) of \(R\)-modules admits a termwise monic
quasi-isomorphism to a K-injective complex with injective terms.
\end{theorem}
\begin{proof}
Put \(T_n=\tau_{\geq -n}X\), where the cutoff term in degree
\(-n\) is \(X^{-n}/d X^{-n-1}\). The maps
\(T_{n+1}\to T_n\) are surjective in each degree, and for
each fixed degree the tower stabilizes at the corresponding
term of \(X\). Thus \(\varprojlim T_n=X\).

We construct compatible termwise monic quasi-isomorphisms
\(T_n\to I_n\) to bounded-below complexes of injectives
whose transitions are split epimorphisms in each degree.
Start with a bounded-below replacement of \(T_0\) by
\cite[R01]{ha-core}. Given \(I_n\), choose a bounded-below
replacement \(T_{n+1}\hookrightarrow J\).
The composite \(T_{n+1}\to T_n\to I_n\) extends to \(g:J\to I_n\)
by Lemma~\ref{res:U04} and bounded-below K-injectivity
\cite[R03]{ha-core}. Put \(C=\mathrm{Cone}(1_{I_n})[-1]\).
With degree-\(i\) term \(I_n^i\oplus I_n^{i-1}\), its differential is
\((a,b)\mapsto(d a,-a-d b)\). Projection onto the first entry is a
chain map \(C\to I_n\) and is surjective with a graded section.
The homotopy \((a,b)\mapsto(-b,0)\) contracts \(C\).
Use \(I_{n+1}=J\oplus C\), map \(T_{n+1}\) into the first summand,
and take transition \(g+\mathrm{pr}_1\).
The map from \(T_{n+1}\) is monic and a quasi-isomorphism,
and commutativity with the transition to \(I_n\) is strict.

Let \(I=\varprojlim I_n\).
Lemma~\ref{res:U05} proves that the natural map \(X\to I\)
is a quasi-isomorphism and that \(I\) is K-injective.
It is monic termwise by compatibility with all the monomorphisms
\(T_n\hookrightarrow I_n\) and \(\varprojlim T_n=X\).
For each degree \(i\), choose splittings of the transition maps
\(I_{n+1}^i\to I_n^i\). Then \(I^i\) is isomorphic as a module to
\(I_0^i\) times the product of their kernels. Each kernel is a direct
summand of an injective module. A product of injectives is injective,
since one can extend a map along a monomorphism in each coordinate.
This proves the final assertion.
\end{proof}

\section{The closed derived structure}

\begin{theorem}[Unbounded derived tensor]\label{res:U07}
There is a symmetric monoidal product
\(\otimes_R^{\mathbf L}:D(R)\times D(R)\to D(R)\) with unit \(R[0]\).
It is computed by replacing both arguments with the models in
Theorem~\ref{res:U02}. If either argument is already K-flat, one
may use the ordinary tensor product without replacing the other.
\end{theorem}
\begin{proof}
By \cite[R04]{ha-core}, K-projective models lift all derived
morphisms uniquely as homotopy classes of chain maps.
Tensoring those classes produces a bifunctor taking values in
\(D(R)\), and the lifted equivalences between models provide its
natural comparison isomorphisms. If \(P,Q\) are K-flat, then
\(P\otimes Q\) is K-flat: associate its tensor with an acyclic
complex and use acyclicity successively. The ordinary chain
associativity, unit and symmetry
\(x\otimes y\mapsto(-1)^{|x||y|}y\otimes x\)
therefore descend, and satisfy the coherence identities because
they already do so on the chain models. Alternatively, after
replacing one of two K-flat complexes by a K-projective model,
the ordinary comparison remains a quasi-isomorphism by the
K-flatness of the other, by Proposition~\ref{res:U03}.
This proves the step needed to iterate the bifunctor independently
of choices. For a K-flat factor tensoring the replacement map in
the other factor is a quasi-isomorphism, again by that proposition.
\end{proof}

\begin{theorem}[Internal derived Hom and adjunction]\label{res:U08}
Write \(\mathbf R\mathrm{Hom}_R(X,Y)\in D(R)\) for internal
derived Hom. It can be computed by
\(\underline{\mathrm{Hom}}_R(P_X,Y)\) with \(P_X\) K-projective,
or by \(\underline{\mathrm{Hom}}_R(X,I_Y)\) with \(I_Y\)
K-injective. Both agree naturally, and
\[
 \mathrm{Hom}_{D(R)}(Z\otimes_R^{\mathbf L}X,Y)
  \cong \mathrm{Hom}_{D(R)}(Z,\mathbf R\mathrm{Hom}_R(X,Y)).
\]
In particular
\(
 H^q\mathbf R\mathrm{Hom}_R(X,Y)
   \cong\mathrm{Hom}_{D(R)}(X,Y[q]).
\)
\end{theorem}
\begin{proof}
Take the two replacements. The expression
\(\underline{\mathrm{Hom}}_R(P_X,I_Y)\) defines a bifunctor
on derived morphisms by \cite[R04]{ha-core} and passage of Hom
through chain homotopy classes. The two single replacement
expressions compare to this one by quasi-isomorphisms by
Proposition~\ref{res:U03}. If \(P\) is K-flat and \(I\) is
K-injective, then \(\underline{\mathrm{Hom}}_R(P,I)\) is
K-injective. Indeed its Hom from an acyclic \(A\) identifies
as a complex with the Hom from \(A\otimes P\) to \(I\), by the
adjunction of Definition~\ref{res:U01}. That complex is acyclic.
For the adjunction take \(P_X,I_Y\) and use the ordinary
chain tensor--Hom adjunction in the homotopy category.
On the left, \(Z\otimes P_X\) computes derived tensor by
Theorem~\ref{res:U07}, and on the right the Hom target is
K-injective by the preceding argument. Both morphism sets
are therefore derived morphism sets by \cite[R04]{ha-core}.
The last assertion follows from the cycle and boundary description
in Definition~\ref{res:U01}, also with the shift of \(Y\).
\end{proof}

\begin{proposition}[Derived extension of scalars]\label{res:U09}
For a homomorphism of commutative rings \(R\to S\), put
\(L(X)=S\otimes_R^{\mathbf L}X\in D(S)\).
Then \(L\) is left adjoint to restriction of scalars, and has the
natural symmetric monoidal comparison
\[
 L(X\otimes_R^{\mathbf L}Y)\ \cong\
 L(X)\otimes_S^{\mathbf L} L(Y).
\]
For successive changes of commutative rings \(R\to S\to T\), the two left extensions
compose to \(T\otimes_R^{\mathbf L}(-)\).
\end{proposition}
\begin{proof}
If \(P\) is K-flat over \(R\), \(S\otimes_R P\) is K-flat over \(S\)
by associativity tested against acyclic \(S\)-complexes.
Similarly extension of a K-projective is K-projective, using the
underived Hom adjunction against an acyclic \(S\)-complex and
exactness of restriction. On a K-projective source the chain Hom
adjunction therefore computes the derived morphisms on both sides.
Choose semifree \(P_X,P_Y\) for the monoidal formula. Both tensors
are computed by the displayed model comparisons, and their
identification is the associativity isomorphism for ordinary
change of rings. It commutes with symmetry, multiplication and the
unit on the chain models, since \(S\) is concentrated in degree zero.
The same base-changed models give transitivity.
\end{proof}

\chapter{Perfect and finite duality}
\label{ch:derived-duality}

All rings in the affine part of this chapter are commutative.
We distinguish internal Hom over a ring from the Hom groups in its
derived category. For a map \(A\to B\) we keep the \(B\)-action on
the right adjoint explicitly.

\section{Perfect complexes}

\begin{theorem}[Evaluation and duals of perfect complexes]\label{res:U10}
Let \(P\in D(R)\) be represented by a bounded complex of finitely
generated projective modules. Define
\(P^\vee=\mathbf R\mathrm{Hom}_R(P,R)\).
Then \(P^\vee\) is perfect, the natural evaluation comparison
\[
 Y\otimes_R^{\mathbf L}P^\vee
     \longrightarrow\mathbf R\mathrm{Hom}_R(P,Y)
\]
is an isomorphism for every \(Y\), and \(P\) is dualizable
with dual \(P^\vee\). Biduality \(P\to(P^\vee)^\vee\) is an
isomorphism.
\end{theorem}
\begin{proof}
Use the bounded finite projective representative for \(P\).
It is K-projective by \cite[R03]{ha-core} and K-flat by
Proposition~\ref{res:U03}. The ordinary dual
\(D=\underline{\mathrm{Hom}}_R(P,R)\) is bounded finite
projective. The comparison in total degree \(q\) sends
\(y\otimes f\) to the map \(x\mapsto y f(x)\) for
\(|y|+|f|=q\). Expanding the differential gives
\(d y\, f-(-1)^{|y|+|f|} y f d_P\), as required.
It is a termwise isomorphism, since each degree of \(P\)
is finitely generated projective and the diagonals involving
\(P\) are finite. The standard evaluation \(D\otimes_R P\to R\)
sends \(f\otimes x\) to \(f(x)\). Applying the isomorphism just
proved with \(Y=P\), the identity of \(P\) corresponds to a
degree-zero cycle in \(P\otimes D\). This defines coevaluation
\(R\to P\otimes D\). In each degree choose a finite dual basis.
The cycle is the sum of the dual basis tensors over the
finitely many degrees. Both triangle identities now follow
by the dual basis expansions in each degree, with evaluation
on the adjacent \(D,P\) factors. In particular the identities
hold already on the complexes. The bidual map is the
degreewise finite projective double dual, with the symmetry
sign \(x\mapsto[f\mapsto(-1)^{|x||f|}f(x)]\).
It is a chain map by the Hom differential and a
termwise isomorphism.
\end{proof}

\begin{proposition}[Projection through a perfect factor and base change]\label{res:U11}
If \(P\) is perfect, the natural comparison
\[
 P\otimes_R^{\mathbf L}\mathbf R\mathrm{Hom}_R(X,Y)
 \longrightarrow
 \mathbf R\mathrm{Hom}_R(X,P\otimes_R^{\mathbf L}Y)
\]
is an isomorphism. Also, for any \(R\to S\) the natural
change-of-rings map
\[
 S\otimes_R^{\mathbf L}\mathbf R\mathrm{Hom}_R(P,Y)
  \longrightarrow
 \mathbf R\mathrm{Hom}_S(S\otimes_R^{\mathbf L}P,
                          S\otimes_R^{\mathbf L}Y)
\]
is an isomorphism.
\end{proposition}
\begin{proof}
For the first assertion test by morphisms from any \(Z\).
Using Theorem~\ref{res:U10}, morphisms from \(Z\) to the
left side are morphisms from \(P^\vee\otimes^{\mathbf L}Z\)
to \(\mathbf R\mathrm{Hom}_R(X,Y)\). The closed adjunction
in Theorem~\ref{res:U08} and associativity identify them with
morphisms from \(P^\vee\otimes^{\mathbf L}(Z\otimes^{\mathbf L}X)\)
to \(Y\). Applying the same duality a second time and
then the closed adjunction gives morphisms from \(Z\) to the
right side. The composite is the canonical tensor-evaluation
comparison, since each step takes adjoints under evaluation
or coevaluation. The isomorphism follows by Yoneda.
For base change, extend scalars on a bounded finite
projective representative for \(P\) and its ordinary dual.
Finite projective duals commute termwise with extension,
and the comparison is compatible with evaluation.
Proposition~\ref{res:U09} transports the evaluation comparison
of Theorem~\ref{res:U10} to the evaluation comparison over \(S\).
This proves the assertion even if \(S\) is not flat over \(R\).
\end{proof}

\begin{example}[Termwise free need not be K-flat]\label{res:U12}
Let \(k\) be a field and \(R=k[\epsilon]/(\epsilon^2)\).
The unbounded complex \(A\) having \(R\) in every degree and
differential multiplication by \(\epsilon\) is acyclic.
It is neither K-projective nor K-flat.
\end{example}
\begin{proof}
Its cycles and boundaries are both \(\epsilon R\) in every degree.
It cannot be K-projective, since its identity is not null-homotopic.
In any proposed homotopy \(d h+h d\) has image in \(\epsilon R\).
For the tensor complex \(A\otimes A\), use in degree \(n\) the
finitely supported basis \(e_i^{(n)}\), indexed by \((i,n-i)\).
Then
\[
 d(e_i^{(n)})=\epsilon e_{i+1}^{(n+1)}
                 +(-1)^i\epsilon e_i^{(n+1)}.
\]
The element \(\epsilon e_0^{(n)}\) is a cycle. If it were a boundary,
a nonzero finitely supported sequence \((a_i)_i\) over \(k\)
would have image \(e_0^{(n)}\) under
\(e_i^{(n-1)}\mapsto e_{i+1}^{(n)}+(-1)^i e_i^{(n)}\).
For such a sequence the least and greatest image indices are its
least index and one more than its greatest index. They are
distinct, a contradiction. This shows directly that \(A\otimes A\)
is not acyclic.
\end{proof}

\begin{example}[An infinite source in Hom base change]\label{res:U13}
Let \(R=k\) be a field, \(S=k[t]\), \(P=\bigoplus_{n\geq0}k\),
and \(Y=k\), all in degree zero. The canonical Hom base-change
map for \(P,Y\) is not surjective.
\end{example}
\begin{proof}
The map is \(k[t]\otimes_k\prod_{n\geq0}k\to\prod_{n\geq0}k[t]\),
by ordinary projectivity and flatness. Any element of its image
has a common bound on the degrees of its entries, since it is a
finite sum of polynomial multiples of scalar sequences.
The sequence \((t^n)_{n\geq0}\) violates that condition.
\end{proof}

\section{Finite affine adjunction and traces}

\begin{theorem}[Derived coinduction and its trace]\label{res:U14}
For a ring map \(A\to B\) let \(f_*:D(B)\to D(A)\) be
restriction of scalars. It has a right adjoint \(f^!\).
If \(Y\to I\) is a K-injective replacement over \(A\), it is computed
by \(f^!Y=\underline{\mathrm{Hom}}_A(B,I)\) with
\((b\cdot\varphi)(b')=\varphi(bb')\).
The counit \(f_*f^!Y\to Y\) is evaluation at \(1\).
Under successive ring maps these right adjoints compose and
the counits compose by successive evaluations at \(1\).
\end{theorem}
\begin{proof}
For any \(B\)-complex \(M\) there is an isomorphism of Hom complexes
\[
 \underline{\mathrm{Hom}}_B
  (M,\underline{\mathrm{Hom}}_A(B,I))
 \longrightarrow \underline{\mathrm{Hom}}_A(f_*M,I),
 \qquad u\longmapsto [m\mapsto u(m)(1)].
\]
Its inverse sends \(v\) to \(m\mapsto[b\mapsto v(bm)]\).
Since \(B\) is in degree zero both are chain maps by the ordinary
Hom differential. Exactness of restriction and K-injectivity
of \(I\) show that \(\underline{\mathrm{Hom}}_A(B,I)\)
is K-injective over \(B\). The displayed isomorphism of degree-zero
cohomologies therefore gives the derived adjunction by
\cite[R04]{ha-core}. It identifies the image of the identity
with evaluation at \(1\), hence gives the counit.
For \(A\to B\to C\) apply the same expression with \(M=C[0]\).
It gives a \(C\)-linear isomorphism
\[
 \underline{\mathrm{Hom}}_B(C,\underline{\mathrm{Hom}}_A(B,I))
 \cong\underline{\mathrm{Hom}}_A(C,I).
\]
The inner model on the left is K-injective, so the left side indeed
computes the composite. Since evaluation at \(1_C\) in the
image means evaluating first at \(1_C\) and then at \(1_B\),
the trace statement follows.
\end{proof}

\begin{proposition}[Finite relative dual under a perfectness hypothesis]\label{res:U15}
If \(B\) is perfect as an \(A\)-complex, put
\(D_{B/A}=f^! A\in D(B)\).
For every \(Y\in D(A)\) the natural \(B\)-linear map
\[
 D_{B/A}\otimes_B^{\mathbf L}(B\otimes_A^{\mathbf L}Y)
     \longrightarrow f^!Y
\]
is an isomorphism.
\end{proposition}
\begin{proof}
Choose a K-injective representative \(I\) of \(A\) and a
semifree representative \(P_Y\) for \(Y\).
The left side is represented by
\(\underline{\mathrm{Hom}}_A(B,I)\otimes_A P_Y\)
because \(B\otimes_A P_Y\) is K-flat over \(B\) by the proof
of Proposition~\ref{res:U09}. Its canonical map to
\(\underline{\mathrm{Hom}}_A(B,I\otimes_A P_Y)\)
sends a pure tensor \(\varphi\otimes p\) to
\([b\mapsto\varphi(b)\otimes p]\).
The corresponding map for a general complex \(Q\) in place of
\(B\) uses \((-1)^{|p||q|}\varphi(q)\otimes p\) on \(q\in Q\).
This is a chain map by the two differentials of
Definition~\ref{res:U01}. Compose with Hom of a
K-injective replacement of \(I\otimes_A P_Y\), which represents
\(Y\). To see this map is a quasi-isomorphism, forget to
\(A\) and resolve \(B\) by a bounded finite projective complex
\(Q\). Replacing \(B\) by \(Q\) in the two Hom complexes with
K-injective targets is a quasi-isomorphism by
Proposition~\ref{res:U03}. Tensoring the comparison for \(I\)
with \(P_Y\) preserves this property. For the \(Q\)-model,
ordinary tensor evaluation is an isomorphism by
Theorem~\ref{res:U10}, and Hom from \(Q\) to the chosen
replacement is a quasi-isomorphism by its K-projectivity.
Commutativity of the comparison squares proves the original map
is a quasi-isomorphism. This construction retains \(B\)-linearity
because the map on the unreplaced \(B\)-models does.
\end{proof}

\begin{proposition}[Flat base change for a finite algebra]\label{res:U16}
Let \(A\) be Noetherian, \(B\) be a finite \(A\)-algebra,
and \(A\to A'\) be flat. Set \(B'=A'\otimes_A B\).
For \(Y\in D^+(A)\), coinduction satisfies the \(B'\)-linear
base-change isomorphism
\[
 B'\otimes_B^{\mathbf L}f^!Y
   \ \cong\ f'^!(A'\otimes_A^{\mathbf L}Y).
\]
It identifies the two evaluation-at-\(1\) traces.
\end{proposition}
\begin{proof}
Take a bounded-below termwise injective representative \(I\) of
\(Y\) and a bounded-below injective replacement \(I'\) of
\(A'\otimes_A I\), using \cite[R01, R03]{ha-core}.
Since \(B\) is finitely presented as an \(A\)-module, flat
extension commutes with \(\mathrm{Hom}_A(B,I^q)\) in each
degree, by taking the kernel of Hom applied to a finite
presentation. Thus the comparison is the ordinary Hom
of \(B'\) with \(A'\otimes_A I\to I'\) on these models.
The \(B\)-base change is underived on every term because
\(B'\) is \(B\)-flat.

For the quasi-isomorphism assertion resolve \(B\) over \(A\)
by a resolution \(P_B\) in nonpositive degrees of finite free
modules. Such a resolution is available by iterating finite
free surjections onto successive finite kernels over the
Noetherian ring. Its scalar extension resolves \(B'\)
by flatness. In each total degree the products in
\(\underline{\mathrm{Hom}}_A(P_B,I)\) are finite because
\(P_B\) is bounded above and \(I\) bounded below. Hence
ordinary Hom on this model commutes with scalar extension
termwise. On the base-changed model, Hom from the
bounded-above projective resolution to \(A'\otimes I\to I'\)
is a quasi-isomorphism by \cite[R03]{ha-core}.
Comparing the two projective models with Hom against
the injective targets proves the claim, since flat tensor
preserves the comparison on the \(A\)-side.
For trace, evaluation at \(1_B\) commutes with change of
scalars and comparison of the target complexes. That is the
counit of Theorem~\ref{res:U14}.
\end{proof}

\begin{corollary}[A quotient by a regular sequence]\label{res:U17}
Let \(A\) be Noetherian, \(J=(a_1,\ldots,a_c)\) be generated
by an \(A\)-regular sequence, and \(B=A/J\).
Then
\[
  f^!Y\ \cong\
  \bigl(\textstyle\bigwedge\nolimits^c_B(J/J^2)\bigr)^\vee[-c]
       \otimes_B^{\mathbf L}(B\otimes_A^{\mathbf L}Y).
\]
Here the dual on the determinant line is its ordinary module dual.
\end{corollary}
\begin{proof}
The Koszul complex on the \(a_i\) resolves \(B\) in degrees
\(-c,\ldots,0\) by \cite[RK-07]{ca-core}. Its graded free
algebra is the exterior algebra on a free module \(F\) with
basis sent to the \(a_i\). In degree \(i\), identify the dual
Koszul term with \(\bigwedge^{c-i}F\otimes(\bigwedge^c F)^\vee\)
by \(z\otimes\lambda\mapsto[x\mapsto\lambda(z\wedge x)]\).
For \(x\in\bigwedge^{i+1}F\) the identity
\(d(z\wedge x)=d z\wedge x+(-1)^{c-i}z\wedge d x=0\)
shows that the Hom differential acts on \(z\) as
\((-1)^{i+1}(-1)^{c-i+1}d z=(-1)^c d z\).
This is the shift by \(-c\). Thus
its cohomology vanishes outside degree \(c\), where it is
\(B\otimes_A(\bigwedge^c F)^\vee\). The action of \(B\) on
this Ext group is the coinduction action. To check it, lift an
element of \(B\) to \(A\). Multiplication on the resolution
by the lift is a chain map lifting multiplication on \(B\),
and induces the ordinary scalar action on cohomology.
If the lift differs by an element of \(J\), the corresponding
multiplication is null-homotopic on the Koszul resolution by
wedge with the corresponding degree \(-1\) element.

The basis of \(F\otimes_A B\) identifies with the classes of
the \(a_i\) in \(J/J^2\). To justify independence, whenever
\(\sum_i b_i a_i\) lies in \(J^2\), subtract a relation with
coefficients in \(J\). The remaining exact relation is in
the image of the degree-two Koszul differential and hence its
coefficients lie in \(J\). Surjectivity follows from generation.
A change of regular generators induces the corresponding
matrix on this conormal module. Lifting its coefficients
gives a comparison of Koszul complexes whose exterior
powers identify the top Ext via the dual determinant of
that same matrix. It follows that the top Ext identification
is canonical. A complex of \(B\)-modules with this unique
nonzero cohomology is canonically represented in \(D(B)\)
by that module in degree \(c\), by good truncation.
The Koszul resolution makes \(B\) perfect over \(A\), so
Proposition~\ref{res:U15} gives the formula for \(Y\).
\end{proof}

\section{Comparison with geometric and local duality}

\begin{proposition}[The finite morphism between affines]\label{res:U18}
Let \(f:\mathrm{Spec}(B)\to\mathrm{Spec}(A)\) be finite.
On the derived categories of complexes of quasi-coherent
modules on these affine schemes, \(f_*\), \(\mathbf Lf^*\),
and the right adjoint to \(f_*\) correspond respectively to
restriction, derived scalar extension, and the coinduction
in Theorem~\ref{res:U14}. The finite trace in this adjunction
corresponds to evaluation at \(1\).
\end{proposition}
\begin{proof}
We recall the needed exact affine correspondence.
For an \(A\)-module \(M\), the module sheaf \(\widetilde M\)
has sections \(M_a\) on \(D(a)\).
Conversely for a quasi-coherent \(\mathcal F\) choose a finite
distinguished affine cover on which it is a module sheaf.
Writing the sheaf equalizer for sections over this cover
and then localizing at \(a\) gives the equalizer for the
induced cover of \(D(a)\), since localization is exact
and the products in this equalizer are finite. Hence
\(\Gamma(X,\mathcal F)_a=\Gamma(D(a),\mathcal F)\).
The canonical map \(\widetilde{\Gamma(X,\mathcal F)}\to
\mathcal F\) is an isomorphism on the distinguished basis.
This equivalence is exact because the module sheaf functor
is exact on each stalk by exactness of localization. It
induces an exact equivalence on derived categories by
termwise application and localization. Preimages of \(D(a)\)
under \(f\) are \(D(a)\) in \(\mathrm{Spec}(B)\), where \(a\)
is viewed in \(B\). Direct image thus corresponds to
restriction of scalars, while ordinary pullback on affines
is extension of scalars by its tensor description.
Apply Proposition~\ref{res:U09} and
Theorem~\ref{res:U14}. The trace is their adjunction counit,
hence the requested evaluation.
\end{proof}

\begin{theorem}[Comparison with complete local duality]\label{res:U19}
Let \((R,\mathfrak m,k)\) be a complete Noetherian local ring.
Choose a surjective local map \(f:S\to R\) from a complete
regular local ring \((S,\mathfrak n,k)\) of dimension \(d\),
with the same residue field. Then
\[
 D_R=f^!(S[d])=\mathbf R\mathrm{Hom}_S(R,S[d])
\]
is the normalized dualizing complex of
\cite[D05]{ca-topics}. For finite \(R\)-modules \(M\)
there are natural isomorphisms
\[
 \begin{split}
 \mathbf R\mathrm{Hom}_R(M,D_R)
      &\cong \mathbf R\mathrm{Hom}_S(M,S[d]),\\
 H^{-q}\mathbf R\mathrm{Hom}_R(M,D_R)
      &\cong \mathrm{Hom}_R(H^q_{\mathfrak m}(M),E_R(k)).
 \end{split}
\]
In the first line the \(R\)-action on the right is by
precomposition on \(M\). The finite trace is the same
evaluation as in Theorem~\ref{res:U14}.
\end{theorem}
\begin{proof}
The regular cover and the identification of the normalized
dualizing complex are \cite[C04, D05]{ca-topics}.
The first isomorphism is the complex-level adjunction of
Theorem~\ref{res:U14}. Multiplication on the \(R\)-factor
there gives multiplication by precomposition on \(M\),
as seen from the inverse formula for \(v\).

Here is a direct comparison with the local cohomology
proof. Put \(E_S=E_S(k)\) and let \(C_{\mathfrak n}\)
be the finite flat augmented \v{C}ech complex for
generators of \(\mathfrak n\). By
\cite[L05, D01, M03]{ca-topics}, \(C_{\mathfrak n}\)
represents \(E_S[-d]\) and
\(\mathrm{End}_S(E_S)=S\).
The complex \(W=\underline{\mathrm{Hom}}_S(C_{\mathfrak n},E_S)\)
is bounded and termwise injective, since Hom from a flat
module into an injective module is injective by tensor--Hom
adjunction and flatness. Exactness of dualization into
\(E_S\) shows that \(W\) has only one nonzero cohomology
module, the copy of \(S\) in degree \(-d\). A cycle representing
its generator \(1\) gives a quasi-isomorphism \(S[d]\to W\).
Thus it may be used for the K-injective model in
Theorem~\ref{res:U14}. Since the images of the generators
generate \(\mathfrak m\), scalar extension identifies
\(R\otimes_S C_{\mathfrak n}=C_{\mathfrak m}\)
by localization. Coinduction identifies
\(E_R(k)\cong\mathrm{Hom}_S(R,E_S)\) by
\cite[D04]{ca-topics}. The complex adjunction yields
\[
 \underline{\mathrm{Hom}}_S(R,W)
    \cong \underline{\mathrm{Hom}}_R(C_{\mathfrak m},E_R).
\]
This is a K-injective \(R\)-model by Theorem~\ref{res:U14}.
Applying \(\underline{\mathrm{Hom}}_R(M,-)\) gives
\(\underline{\mathrm{Hom}}_R(M\otimes_R C_{\mathfrak m},E_R)\).
Taking degree \(-q\) cohomology gives the displayed Matlis
dual, since Hom into an injective module is exact and the
tensor computes local cohomology. The comparison agrees
with \cite[D06]{ca-topics}, using the same bounded
injective model and the same change of scalars. The trace
on the coinduced model is evaluation at \(1_R\) as claimed.
\end{proof}

\begin{corollary}[Completion and the Cohen--Macaulay form]\label{res:U20}
For a Noetherian local ring \(R\) and a finite \(R\)-module \(M\)
there is an isomorphism
\[
 \mathrm{Hom}_R(H^q_{\mathfrak m}(M),E_R)
   \cong H^{-q}\mathbf R\mathrm{Hom}_{\widehat R}
                             (\widehat M,D_{\widehat R})
\]
of \(\widehat R\)-modules. If \(R\) is complete Cohen--Macaulay
of dimension \(e\), this becomes
\(\mathrm{Ext}^{e-q}_R(M,\omega_R)\), where
\(\omega_R=H^{-e}(D_R)\).
\end{corollary}
\begin{proof}
The local cohomology groups are \(\mathfrak m\)-power torsion
and are unchanged on completing \(M\), by \cite[L02, L09]{ca-topics}.
Every \(R\)-linear map between \(\mathfrak m\)-power torsion
modules is \(\widehat R\)-linear. This follows by using
\(R/\mathfrak m^r=\widehat R/(\mathfrak m\widehat R)^r\)
on submodules annihilated by \(\mathfrak m^r\). The
\(\widehat R\)-injective envelope of \(k\) restricts
to the \(R\)-injective envelope, as in \cite[M06]{ca-topics}.
Thus the first isomorphism follows by applying
Theorem~\ref{res:U19} to \(\widehat R\). For the second,
Theorem~\ref{res:U19} with \(M=R\) and the
Cohen--Macaulay cohomology criterion \cite[V07]{ca-topics}
show that \(D_R\) has a single nonzero cohomology
module in degree \(-e\). The module is the canonical
module by \cite[D08]{ca-topics}, and good truncation
identifies \(D_R\) with \(\omega_R[e]\). Now the
cohomology shift gives the claimed Ext group. These
are the completions and shifts of \cite[D07--D09]{ca-topics}.
\end{proof}

\section{An explicit projective case}

\begin{theorem}[Laurent monomials compute the twists on projective space]\label{res:U21}
Let \(k\) be a field, \(n\geq1\), \(X=\mathbf P^n_k\),
and \(S=k[x_0,\ldots,x_n]\) with its usual grading.
For \(m\in\mathbb Z\), the cohomology of \(\mathcal O_X(m)\)
is \(S_m\) in degree zero, vanishes for \(0<q<n\)
and for \(q>n\), and has basis
\[
   \{x_0^{a_0}\cdots x_n^{a_n}:\ a_i<0,\
           a_0+\cdots+a_n=m\}
\]
in degree \(n\).
\end{theorem}
\begin{proof}
Use the cover \(D_+(x_i)\). Its nonempty intersections
are affine Noetherian schemes, since the one-variable
dehomogenizations on a chart are polynomial rings and
localizations. Sections of \(\mathcal O(m)\) on the
intersection indexed by a nonempty subset \(J\) are the
degree-\(m\) part of \(S[x_j^{-1}:j\in J]\), by the
construction of the twisting sheaf. Quasi-coherent
modules on these affine intersections have vanishing
positive cohomology by the affine acyclicity argument of
\cite[L10]{ca-topics}. The \v{C}ech comparison
\cite[H02--H04]{ha-core} thus computes the cohomology.

In the \v{C}ech complex, decompose into exponent vectors
\(\mathbf a=(a_0,\ldots,a_n)\) of total degree \(m\).
The term for a subset \(J\) has a basis vector for
\(\mathbf a\) precisely when it contains the set
\(N=\{i:a_i<0\}\). Restrictions either carry the basis
vector to itself or have zero source, so each vector
gives a subcomplex. When \(N\) is nonempty and proper,
choose \(j\notin N\). The alternating simplex homotopy
pairs the components \(J\) not containing \(j\) with
\(J\cup\{j\}\). Explicitly the degree-minus-one map
takes the component at \(J\cup\{j\}\) to the component
at \(J\) with sign \((-1)^{|\{i\in J:i<j\}|}\), and has
zero component at \(J\) if \(j\in J\).
Both paired sets are admissible, because \(J\) contains
\(N\) and is therefore nonempty. Expanding the
alternating differential gives the identity on the
insertion of \(j\), while insertions of indices other
than \(j\) occur twice with opposite signs. Thus this
summand is acyclic. For \(N=\varnothing\), the same
homotopy works after adding the empty subset in
degree \(-1\), leaving exactly one copy of \(k\) in
degree zero of the unaugmented subcomplex.
For \(N\) the full set only the top-degree term remains.
Degreewise direct sums over exponent vectors preserve
cohomology. This proves the formulas.
\end{proof}

\begin{corollary}[Trace and duality for the twists]\label{res:U22}
In the situation of Theorem~\ref{res:U21}, the coefficient
of \((x_0\cdots x_n)^{-1}\) specifies a trace
\(H^n(X,\mathcal O_X(-n-1))\to k\).
Multiplication followed by trace gives perfect pairings
\[
 H^q(X,\mathcal O_X(m))\ \otimes_k\
 H^{n-q}(X,\mathcal O_X(-m-n-1))\ \longrightarrow k
 \qquad (0\leq q\leq n).
\]
The sheaf \(\mathcal O_X(-n-1)\) identifies with the
top differential forms \(\Omega^n_{X/k}\).
\end{corollary}
\begin{proof}
The trace space is one-dimensional by
Theorem~\ref{res:U21}. The only potentially nonzero
pairings are \(q=0,n\). In each of them multiplication
by a global section acts on the Laurent \v{C}ech
representative by ordinary multiplication. If \(x^{\mathbf a}\)
has all exponents nonnegative, its unique complementary
basis vector has exponents \(-1-a_i\) at each coordinate.
Their product gives coefficient \(1\), while any other
basis product gives coefficient zero. Thus the bases in
the two factors are dual. The formula also covers degrees
for which both spaces are zero.

For the top forms, on \(U_i=D_+(x_i)\) write
\(t_a=x_a/x_i\) and use the generator
\[
 \eta_i=(-1)^i\,d t_0\wedge\cdots
                   \wedge\widehat{d t_i}\wedge\cdots\wedge d t_n .
\]
For \(j\ne i\) put \(z=t_j=x_j/x_i\), a unit on the overlap.
The new coordinate indexed by \(i\) is \(1/z\), with
differential \(-z^{-2}dz\). Every other coordinate indexed
by \(a\ne i,j\) is \(t_a/z\), with differential
\(z^{-1}d t_a-z^{-2}t_a\,dz\). In wedging the new
differentials the summands with a second \(dz\) vanish.
The scalar factor without its sign is \(z^{-n-1}\), and
the sign of \(-d z\) moved from the \(i\)-slot to the \(j\)-slot
is \((-1)^{i+j}\). The factors \((-1)^i,(-1)^j\)
in the two generators cancel that sign. Hence
\(\eta_j=z^{-n-1}\eta_i\), the transition function of
the frame \(x_i^{-n-1}\) for \(\mathcal O(-n-1)\).
This gives the asserted identification.
\end{proof}

\chapter{Model Categories and Projective Complexes}
\label{ch:model-categories}

A model category organizes localization by giving two compatible
procedures for changing an object or a map. We develop the homotopy
category and the Quillen adjunction directly from lifting, then establish
the unbounded projective model on complexes. Throughout, maps of
complexes have degree zero, and cochain differentials raise degree.
The conventions for model categories are those of \cite[Chapter 1]{hovey}
and \cite[Section 11.3]{riehl}.

\section{Lifting and homotopy}

For two maps \(i:A\to B\) and \(p:X\to Y\), write \(i\perp p\)
when every commutative square with \(i\) on the left and \(p\) on
the right has a diagonal filler \(B\to X\). A retract of arrows is
a retract in the category of commutative squares.

\begin{definition}[Model structure]\label{res:Q01}
A model category in this chapter is a complete and cocomplete category
with classes \(W\), \(\operatorname{Cof}\), and \(\operatorname{Fib}\)
of weak equivalences, cofibrations, and fibrations. The classes contain
isomorphisms and are closed under retracts. The class \(W\) satisfies
two-out-of-three. Call a map \emph{trivial} when it belongs to \(W\).
Cofibrations lift against trivial fibrations, and trivial cofibrations
lift against fibrations. Each map has functorial factorizations as a
cofibration followed by a trivial fibration and as a trivial cofibration
followed by a fibration.
Write \(0\) and \(1\) for the initial and terminal objects. An object
\(A\) is cofibrant if \(0\to A\) is a cofibration and fibrant if
\(A\to1\) is a fibration.
\end{definition}

The two lifting pairs determine their classes by orthogonality. Indeed,
if \(h\perp\operatorname{Fib}\), factor \(h=p i\) with \(i\) a
trivial cofibration and \(p\) a fibration. The lift against \(p\)
exhibits \(h\) as a retract of \(i\). The argument for the other
left class uses the other factorization, and duality gives the
right classes. A left class so determined is closed under pushout,
composition, transfinite composition, and retract. For a pushout,
the lift for the old map extends by the defining pushout. For an
ordinal sequence choose lifts successively and use the defining
colimit at limits. A right class has the dual closures, in
particular pullback and composition. These arguments supply the
closure properties used below.

\begin{proposition}[Cylinder and path objects]\label{res:Q02}
A cylinder for \(A\) is a factorization
\[
 A\amalg A\xrightarrow{(j_0,j_1)} C\xrightarrow{r}A
\]
of the fold map with first map a cofibration and \(r\in W\).
It is called good when \(r\) is a trivial fibration.
A path object for \(B\) is a factorization
\[
 B\xrightarrow{s} P\xrightarrow{(e_0,e_1)}B\times B
\]
of the diagonal with \(s\in W\) and second map a fibration.
It is called good when \(s\) is a trivial cofibration.
Good choices exist functorially. If \(A\) is cofibrant, the
two endpoints of any cylinder are trivial cofibrations. If \(B\)
is fibrant, both evaluations of any path object are trivial
fibrations.
\end{proposition}
\begin{proof}
Apply the two factorizations of \ref{res:Q01}. For cofibrant \(A\)
each coproduct inclusion into \(A\amalg A\) is a cofibration by
pushout from \(0\to A\). Therefore \(j_i\) is a cofibration,
and \(rj_i=\mathrm{id}\) gives \(j_i\in W\). For fibrant \(B\),
each product projection out of \(B\times B\) is a fibration
by pullback from \(B\to1\). Now \(e_i s=\mathrm{id}\) gives
\(e_i\in W\).
\end{proof}

\begin{theorem}[The homotopy relation]\label{res:Q03}
For cofibrant \(A\) and fibrant \(B\), left homotopy and right
homotopy are the same equivalence relation on maps \(A\to B\).
Left homotopy from \(f\) to \(g\) is a map \(C\to B\) extending
\((f,g)\) on the two ends of a cylinder. Right homotopy is a
map \(A\to P\) with evaluations \(f,g\). Any fixed cylinder
or path object may be used. Homotopy is preserved by
precomposition from a cofibrant object and postcomposition
into a fibrant object. If \(f\) and \(g\) are homotopic,
then \(f\in W\) if and only if \(g\in W\).
\end{theorem}
\begin{proof}
Suppose a left homotopy \(H:C\to B\) is provided. Factor \(r\)
as a trivial cofibration \(C\to C'\) followed by a fibration
\(C'\to A\), which is trivial by two-out-of-three. The first map
allows \(H\) to extend, since \(B\) is fibrant. This gives a
homotopy on a good cylinder. A map from any fixed cylinder to
that good cylinder, preserving both ends and commuting with the
collapse to \(A\), is obtained by lifting. Thus every fixed
cylinder realizes the same left relation. Dually every fixed
path object realizes the same right relation if \(A\) is
cofibrant.

For the equivalence of the relations, choose a good cylinder
and a good path object. A left homotopy \(H\) gives a lift
\(C\to P\) against \(P\to B\times B\) with top map \(s f\)
on the \(j_0\) end and bottom map \((fr,H)\).
Its restriction to \(j_1\) is a right homotopy. Conversely
if \(K:A\to P\) is a right homotopy, lift
\((s f,K):A\amalg A\to P\) against \(e_0:P\to B\)
with bottom map \(fr:C\to B\). Composing the lift with
\(e_1\) gives the left homotopy.

Reflexivity uses \(fr\) and symmetry interchanges the ends.
For transitivity glue two cylinders along the second end
of the first and first end of the second. The outer end
inclusion is a cofibration. It factors through the first
cylinder together with one extra copy of \(A\), after which
one pushes out the boundary inclusion of the second cylinder.
The inclusion of the first cylinder into the glued object
is a pushout of a trivial endpoint cofibration. Hence the
common collapse of the glued object is in \(W\) by
two-out-of-three. The two homotopies agree on the middle end,
so they glue to a homotopy on this cylinder.

Postcomposition of left homotopies is literal. For
precomposition along \(A'\to A\) with \(A'\) cofibrant,
use a good cylinder on \(A\) and lift from a cylinder on
\(A'\), preserving the two new end maps. Finally the
endpoints are both weak equivalences by \ref{res:Q02}.
Apply two-out-of-three to their composites with \(H\).
\end{proof}

\begin{proposition}[Whitehead criterion]\label{res:Q04}
A map between cofibrant fibrant objects is a weak equivalence
if and only if it admits a homotopy inverse.
\end{proposition}
\begin{proof}
For a trivial cofibration \(i:A\to Z\) between such objects,
lift against \(A\to1\) to obtain a retraction \(u:Z\to A\).
Lift along \(i\) into a path object of \(Z\) with endpoints
\((iu,\mathrm{id}_Z)\), using the constant path on \(i\).
This gives \(iu\simeq\mathrm{id}_Z\). Dually a trivial
fibration \(p:Z\to B\) between such objects has a section
which is a homotopy inverse. Factoring an arbitrary weak
equivalence as a trivial cofibration followed by a fibration
gives a trivial fibration by two-out-of-three, with
intermediate object again cofibrant fibrant.

Conversely factor a map having a homotopy inverse as a
cofibration \(i:A\to Z\) followed by a trivial
fibration \(p:Z\to B\). The intermediate object is again
cofibrant fibrant. As \(p\) has a homotopy inverse, so does
\(i\), say \(u_0\). Choose a right homotopy from
\(u_0 i\) to \(\mathrm{id}_A\). Lift it along \(i\) against
the first evaluation from a path object of \(A\), with
bottom map \(u_0\). Its other evaluation gives
\(u:Z\to A\) with \(ui=\mathrm{id}_A\) and \(u\simeq u_0\).
Thus \(iu\simeq\mathrm{id}_Z\), and \(iu\in W\) by
\ref{res:Q03}. The map \(i\) is a retract of the arrow \(iu\)
via the pairs \((i,\mathrm{id}_Z)\) and
\((u,\mathrm{id}_Z)\). Therefore \(i\), and hence \(pi\),
belongs to \(W\).
\end{proof}

\begin{theorem}[Localization by homotopy classes]\label{res:Q05}
The category of cofibrant fibrant objects with morphisms the
homotopy classes is a localization of the model category at
\(W\). A map of the model category becomes an isomorphism
in this localization exactly when it belongs to \(W\).
\end{theorem}
\begin{proof}
Composition is well-defined by \ref{res:Q03}. Choose a
functorial cofibrant replacement \(QX\to X\), a trivial
fibration, and a functorial fibrant replacement \(X\to RX\),
a trivial cofibration. The object \(RQX\) is cofibrant
fibrant. Naturality and two-out-of-three show that the
replacement functors preserve \(W\) and reflect it on
the maps in the two natural replacement squares.
The functor \(X\mapsto RQX\) followed by passage to homotopy
classes inverts \(W\) by \ref{res:Q04}.

For a functor \(F\) inverting \(W\), its restriction to the
cofibrant fibrant objects passes through homotopy classes.
Indeed a cylinder collapse goes to an isomorphism and its two
sections consequently go to the same map. The resulting
factorization of \(F\) is natural up to the natural
isomorphisms induced by \(QX\to X\) and \(QX\to RQX\).
For \(X\) already cofibrant fibrant, \(QX\) is also fibrant
because \(QX\to X\) is a fibration. The two replacement maps
therefore identify, naturally in the homotopy category,
the original cofibrant fibrant object and its new replacement.
Restriction and replacement are thus inverse up to natural
isomorphism on functors that invert \(W\), including natural
transformations. This proves the localization property.
If the image of \(f\) is invertible, the representative
\(RQf\) is a homotopy equivalence, so is in \(W\) by
\ref{res:Q04}. Two-out-of-three in the replacement squares
gives \(f\in W\).
\end{proof}

\begin{proposition}[Cofibrant source and fibrant target]\label{res:Q06}
For cofibrant \(A\) and fibrant \(B\), homotopy classes of
actual maps \(A\to B\) map bijectively to morphisms
\(\operatorname{Ho}(\mathcal C)(A,B)\).
\end{proposition}
\begin{proof}
Replace \(A\) by \(A\hookrightarrow RA\), a trivial
cofibration into a cofibrant fibrant object. Every map to
\(B\) extends. If the restrictions of two extensions are
right homotopic, extend the homotopy along \(A\to RA\)
by lifting against \(P\to B\times B\), with the two
extensions as endpoints. Thus this replacement gives a
bijection on homotopy classes. Replace \(B\) by a trivial
fibration \(QB\to B\), so that \(QB\) is cofibrant
fibrant. For source \(RA\), maps lift against this
trivial fibration. If two composites are left homotopic,
lift the cylinder homotopy with the two maps to \(QB\)
as endpoints. Therefore this replacement also gives a
bijection on classes. Both replacements induce
isomorphisms in the localization of \ref{res:Q05}.
This identifies the stated map with a composite of
bijections. In particular the identification is natural
on maps of cofibrant sources and fibrant targets.
\end{proof}

\section{Quillen adjunctions}

\begin{lemma}[Brown's lemma on cofibrant objects]\label{res:Q07}
A functor on the cofibrant objects which carries trivial
cofibrations to weak equivalences carries every weak
equivalence there to a weak equivalence, assuming two-out-of-three
in the target. The corresponding assertion for fibrant
objects and trivial fibrations holds dually.
\end{lemma}
\begin{proof}
Let \(f:X\to Y\) be a weak equivalence of cofibrant objects.
Factor \((f,\mathrm{id}_Y):X\amalg Y\to Y\) as a
cofibration followed by a trivial fibration \(p:Z\to Y\).
Both maps \(X\to Z\) and \(Y\to Z\) are trivial
cofibrations by two-out-of-three, and \(Z\) is cofibrant.
The image of \(p\) is weak because its composite with
the image of \(Y\to Z\) is the identity. The same then
holds for the image of \(f\).
\end{proof}

An adjunction \(L:\mathcal C\rightleftarrows\mathcal D:U\)
is Quillen when \(L\) preserves cofibrations and trivial
cofibrations, equivalently when \(U\) preserves fibrations
and trivial fibrations. This equivalence follows by
transporting lifting squares along the adjunction and
using the orthogonality observation after \ref{res:Q01}.

\begin{theorem}[Derived Quillen adjunction]\label{res:Q08}
A Quillen adjunction induces an adjunction
\[
 \mathbf L L:\operatorname{Ho}(\mathcal C)
 \rightleftarrows \operatorname{Ho}(\mathcal D):\mathbf R U.
\]
It is computed by \(LQ\) and \(UR\), where \(Q\) is
cofibrant replacement in \(\mathcal C\) and \(R\)
is fibrant replacement in \(\mathcal D\).
\end{theorem}
\begin{proof}
The two composites preserve weak equivalences by
\ref{res:Q07} and hence descend by \ref{res:Q05}.
The functors \(L\) and \(U\) preserve respectively
cofibrancy and fibrancy because they preserve the relevant
structure maps and initial or terminal objects.
For cofibrant \(X\) and fibrant \(Y\), the ordinary adjunction
identifies maps \(LX\to Y\) with maps \(X\to UY\).
Apply \(L\) to a cylinder of \(X\). Its end inclusion
is a cofibration and the collapse is a weak equivalence
by \ref{res:Q07}, since both ends of the collapse map
are cofibrant. It is a cylinder of \(LX\) because \(L\)
preserves the coproduct. The ordinary adjunction identifies
maps out of this cylinder with maps out of the original
cylinder on the other side and respects endpoint maps.
By \ref{res:Q03}, the adjunction therefore identifies
the two homotopy relations. Applying \ref{res:Q06} gives
the adjunction on homotopy categories. Replacing both
variables supplies this bijection for arbitrary objects,
naturally by the functorial choices and the natural
identification of \ref{res:Q06}.
\end{proof}

\begin{theorem}[Quillen equivalence criterion]\label{res:Q09}
The derived adjunction of \ref{res:Q08} is an equivalence
if and only if the following condition holds. For every
cofibrant \(X\) and fibrant \(Y\), a map \(LX\to Y\) is
a weak equivalence exactly when its adjoint \(X\to UY\)
is a weak equivalence.
\end{theorem}
\begin{proof}
Assume the condition. For cofibrant \(X\), choose a
fibrant replacement \(LX\to R LX\). Its adjoint is
weak by the condition, and represents the derived unit.
For fibrant \(Y\), choose a cofibrant replacement
\(Q UY\to UY\). Its adjoint is weak and represents
the derived counit. Thus the unit and counit are
isomorphisms by \ref{res:Q05}, and functorial
replacement gives this conclusion on all objects.
Conversely, an adjunction which is an equivalence
identifies isomorphisms with isomorphisms under the
adjunction bijection. Apply \ref{res:Q05} to both
maps in the claimed criterion. This proof uses the
unit at a fibrant replacement of \(LX\) and the
counit at a cofibrant replacement of \(UY\), in agreement
with \cite[correction to Corollary 1.3.16]{hovey-errata}.
\end{proof}


\section{Cochain complexes over a ring}

Let \(A\) be a ring with identity. Write \(\operatorname{Ch}(A)\)
for the category of unbounded cochain complexes of left \(A\)-modules.
It has all small limits and colimits, obtained degree by degree.
We use quasi-isomorphisms as weak equivalences and reindex the
projective disk generators of \cite[Section 2.3]{hovey} cohomologically.

\begin{proposition}[Spheres, disks, and their lifts]\label{res:Q10}
Let \(S^n\) be \(A\) in degree \(n\) and \(D^n\) be
\(A\xrightarrow{\mathrm{id}} A\) in degrees \(n,n+1\).
Use the sets of maps
\[
 I=\{S^{n+1}\longrightarrow D^n:n\in\mathbb Z\},
 \qquad
 J=\{0\longrightarrow D^n:n\in\mathbb Z\},
\]
where \(S^{n+1}\to D^n\) includes the upper degree. A map has
the right lifting property against \(J\) precisely when it is
degreewise surjective. It has the right lifting property against
\(I\) precisely when it is a degreewise surjective quasi-isomorphism.
\end{proposition}
\begin{proof}
A map out of \(D^n\) specifies any element in degree \(n\).
Lifting against \(J\) thus tests degreewise surjectivity.
A square against \(S^{n+1}\to D^n\) for \(p:X\to Y\) specifies
\(z\in Z^{n+1}(X)\) and \(y\in Y^n\) such that \(p(z)=dy\).
A lift is \(x\in X^n\) with \(p(x)=y\) and \(dx=z\).
Taking \(z=0\) gives surjectivity on cycles in degree \(n\).
For general \(y\), first lift the cycle \(dy\) and then
lift \(y\), obtaining degreewise surjectivity. Taking \(y=0\)
shows that the kernel has zero cohomology.

Conversely suppose \(p\) is degreewise surjective and a
quasi-isomorphism. Its kernel is acyclic by the long exact
sequence of cohomology. Start with a graded lift \(x_0\) of
\(y\). The difference \(z-dx_0\) is a cycle in the kernel,
hence the differential of some \(u\) in the kernel. Then
\(x=x_0+u\) gives the required lift.
\end{proof}

\begin{lemma}[Free cell factorizations]\label{res:Q11}
Every cochain map factors functorially both as an \(I\)-cell
map followed by an \(I\)-right lifting map and as a \(J\)-cell
map followed by a \(J\)-right lifting map. Here a cell is
attached by pushing out a coproduct of generators, and a
cell map is a sequential composite of such attachments.
An \(I\)-cell map is obtained by adjoining free graded
generators with differential in the previous stage.
A \(J\)-cell map adjoins a direct sum of disks.
\end{lemma}
\begin{proof}
Start with the given arrow \(X_0\to Y\). Given \(X_t\to Y\),
take the coproduct of one copy of each generator for every
commutative square from that generator to \(X_t\to Y\).
Push out its domain map to obtain \(X_{t+1}\), with the
induced map to \(Y\). In degree \(n\), an \(I\)-attachment
adjoins one free generator \(e\) of degree \(n\) for each
corresponding square, imposing \(de=z\) where \(z\) is the
image of the sphere generator. This is consistent because
\(z\) is closed. For \(J\), the whole disk is added freely
as a direct summand of complexes. Take the sequential
colimit \(X_\infty\). The domain of an \(I\)-generator is
one cycle generator. A map from it to \(X_\infty\) is
represented at a finite stage because the cycle condition
becomes true at some finite stage. A \(J\)-domain is zero.
Thus any square at the colimit factors through a finite
stage, after which it acquires its filler. Using all
squares at each stage makes the factorizations functorial.
\end{proof}

\begin{theorem}[The unbounded projective model structure]\label{res:Q12}
There is a model structure on \(\operatorname{Ch}(A)\)
with weak equivalences the quasi-isomorphisms, fibrations
the degreewise surjections, and cofibrations the maps
lifting against degreewise surjective quasi-isomorphisms.
Every object is fibrant. Its homotopy category is
\(\mathrm D(A)\).
\end{theorem}
\begin{proof}
The weak equivalences have two-out-of-three and are
retract-closed by functorial cohomology. The lifting-defined
classes are retract-closed. By \ref{res:Q10}, the two
factorizations of \ref{res:Q11} have the required right
classes. Every \(I\)-cell map is a cofibration by its
lifting property. Every \(J\)-generator is an \(I\)-cell
map, since \(0\to S^{n+1}\) is the pushout of
\(S^{n+2}\to D^{n+1}\) along the zero map and
\(S^{n+1}\to D^n\) is itself a generator. Consequently
\(J\)-cell maps are cofibrations. They are also
quasi-isomorphisms since adjoining a direct sum of disks
does not change cohomology.

It remains to identify the left lifting class against
fibrations. Every map in this class is a retract of its
\(J\)-cell factor by the lift against the right factor.
It is therefore a trivial cofibration. Conversely
factor a trivial cofibration \(i\) as \(i=pj\) with \(j\)
a \(J\)-cell map and \(p\) a fibration. As \(i\) and \(j\)
are quasi-isomorphisms, \(p\) is a trivial fibration.
The lift furnished by \(i\) exhibits \(i\) as a retract
of \(j\), which has the left lifting property against
every fibration. This verifies both lifting classes and
all the axioms of \ref{res:Q01}. All maps to the zero complex
are surjective. The final assertion follows from \ref{res:Q05}
and the definition of the derived category by localization.
\end{proof}

\begin{proposition}[Cofibrations and acyclic targets]\label{res:Q13}
A map of complexes is a cofibration in \ref{res:Q12}
if and only if it is a degreewise split injection whose
cokernel \(P\) is degreewise projective and \(K\)-projective.
Here \(K\)-projective means that
\(\operatorname{Hom}_A^\bullet(P,E)\) is acyclic for every
acyclic \(E\). The cokernel is also \(K\)-flat, meaning
that \(E\otimes_A P\) is acyclic for every acyclic complex
of right \(A\)-modules.
\end{proposition}
\begin{proof}
A cofibration is a retract of its \(I\)-cell factor.
An \(I\)-cell map is split injective as a graded map,
and its cokernel has a filtration with graded-split
inclusions and freely adjoined generators whose
differentials land in earlier stages. A cochain map \(f\)
from such a cokernel to an acyclic \(E\) is null homotopic.
At the stage adjoining \(e\), with \(de\) already present,
the equation \(f(e)=d h(e)+h(de)\) is solvable because
\(f(e)-h(de)\) is a cycle. Indeed
\(d f(e)-d h(de)=f(de)-f(de)+h(d^2e)=0\)
using the equation on the prior stage. Choose these values
on all generators and extend linearly. The same argument
applied to all shifts of \(E\) is precisely acyclicity
of the Hom complex. The cokernel is \(K\)-flat because
tensoring the graded-split filtration with \(E\) gives
a filtration by complexes with acyclic quotients, direct
sums of shifts of \(E\). Each finite stage is acyclic
by the long exact sequence and the union is acyclic
by exactness of filtered colimits of modules. Taking
retracts preserves all these conditions, which proves
necessity.

For sufficiency let \(i:A_0\to B_0\) have the stated
properties and consider a lifting square against a
surjective quasi-isomorphism \(q:X\to Y\). Using the
graded splitting of \(i\) and degreewise projectivity
of \(P=B_0/A_0\), choose a graded lift \(u:B_0\to X\)
extending the prescribed map on \(A_0\).
Its defect \(d_Xu-u d_{B_0}\) vanishes on \(A_0\)
and has image in \(\ker q\). It defines a degree-one
cocycle of \(\operatorname{Hom}_A^\bullet(P,\ker q)\).
The identity on two successive differentials gives the
cocycle equation. By \ref{res:Q10}, the kernel is
acyclic. Thus the defect is the Hom differential of
a degree-zero map \(P\to\ker q\). Subtract this
map after quotient and inclusion to correct \(u\)
to a cochain map and complete the lift.
\end{proof}

\begin{corollary}[Chain homotopy in the projective model]\label{res:Q14}
For cofibrant \(X\), the morphisms from \(X\) to any
complex \(Y\) in \(\mathrm D(A)\) are cochain maps modulo
ordinary cochain homotopy.
\end{corollary}
\begin{proof}
Use the path complex
\[
 P(Y)^n=Y^n\oplus Y^n\oplus Y^{n-1},\qquad
 d(y_0,y_1,b)=(dy_0,dy_1,y_1-y_0-db).
\]
Its third component after two applications of \(d\) is
zero, and its evaluations are the two projections. They
together form a degreewise surjection to \(Y\times Y\).
The first evaluation has contractible kernel, since in
coordinates \((a,b)\) that kernel has differential
\((da,a-db)\) and contracting homotopy \((a,b)\mapsto(b,0)\).
The diagonal \(y\mapsto(y,y,0)\) is therefore a
quasi-isomorphism. A cochain map \(X\to P(Y)\) with
first two coordinates \(f,g\) has third coordinate
\(h\) with \(g-f=d h+h d\). Now apply \ref{res:Q03}
and \ref{res:Q06}, using fibrancy of all complexes.
\end{proof}

\section{Tensor products and homotopy limits}

\begin{theorem}[Tensor compatibility]\label{res:Q15}
If \(A\) is commutative, the projective model structure
is a closed symmetric monoidal model structure for the
usual tensor product of cochain complexes. In particular,
the pushout product of two cofibrations is a cofibration
and is trivial if either is trivial.
\end{theorem}
\begin{proof}
The tensor product uses direct sums in total degree and
the signed differential. Its internal Hom uses products
in each degree with differential
\(d(h)=d h-(-1)^{|h|}h d\). Evaluation and currying give
the adjunction by expansion of the cochain identities.
Symmetry is \(x\otimes y\mapsto(-1)^{|x||y|}y\otimes x\).

For \(a:A_0\to B_0\), \(b:C_0\to D_0\), their
pushout product is
\[
 (B_0\otimes C_0)\amalg_{A_0\otimes C_0}(A_0\otimes D_0)
 \longrightarrow B_0\otimes D_0.
\]
For \(a,b\in I\), the inclusion has the product of the
two lower disk generators as its only new generator.
Its differential lies in the source, so this is a
single \(I\)-attachment. For \(a\in I\), \(b\in J\),
the inclusion splits degreewise and its quotient is
a shift of a disk, from the lower sphere of the
\(I\)-cokernel tensored with the disk of \(b\).
The sign is absorbed by changing the sign of one disk
generator. The quotient is acyclic, degreewise free
and \(K\)-projective by contractibility. Thus the
pushout product is a trivial cofibration by
\ref{res:Q13}. Symmetry handles \(I,J\) in the other order.

For fixed \(a\), the operation \(a\square -\)
preserves pushouts of arrows by separate cocontinuity
of the tensor product. For a composite \(C\to D\to E\)
in the other variable, first pass through
\[
 (B_0\otimes D)\amalg_{A_0\otimes D}(A_0\otimes E)
\]
on the way to \(B_0\otimes E\). The first map in
this factorization is a pushout of the pushout product
for \(C\to D\), and the second is the pushout product
for \(D\to E\). For a sequence of attachments use these
same intermediates with its final object \(E\), taking
colimits at limit stages. This proves closure under
the cell operations in either variable. Retracts are
preserved because the pushout product is functorial
on pairs of arrows. Since cofibrations are retracts
of \(I\)-cell maps and trivial cofibrations retracts
of \(J\)-cell maps by \ref{res:Q12}, the generator
calculations prove the assertion. The unit \(A[0]\)
is cofibrant by \ref{res:Q13}, so the identity unit
replacement verifies the unit axiom. For any weak
cofibrant replacement of this unit the same conclusion
follows from \ref{res:Q07} applied to tensor with a
cofibrant object, which is left Quillen by the
pushout-product statement.
\end{proof}

Use the cone convention
\(\operatorname{Cone}(h)^n=Z^n\oplus T^{n+1}\)
for \(h:T\to Z\), with differential
\(d(z,t)=(dz+h(t),-dt)\).
The cochain shift \([-1]\) negates the differential.
For the general homotopy-limit viewpoint see
\cite[Chapters 5--6]{riehl}. The explicit models below use exact
products in module categories, as in \cite[Section 34]{st-derived}.

\begin{proposition}[Homotopy pullbacks of complexes]\label{res:Q16}
For a span \(X\xrightarrow{f} Z\xleftarrow{g}Y\),
a homotopy pullback is
\[
 \operatorname{Cone}(f-g:X\oplus Y\longrightarrow Z)[-1].
\]
It is functorial on strict maps of spans and invariant under
objectwise quasi-isomorphism. When either leg is a fibration,
the strict pullback maps to this model by a quasi-isomorphism.
\end{proposition}
\begin{proof}
Put \(C^n=Z^n\oplus Z^{n-1}\) with differential
\((a,b)\mapsto(da,a-db)\). The homotopy
\((a,b)\mapsto(b,0)\) contracts \(C\). Replace \(f\)
by the surjection \(X\oplus C\to Z\) given by
\((x,a,b)\mapsto f(x)+a\). The original \(X\) includes
by a quasi-isomorphism over \(Z\). The strict pullback
of the resulting span consists of triples
\((x,y,b)\) with \(a=g(y)-f(x)\), and its differential
is \((dx,dy,g(y)-f(x)-db)\). This is precisely the
stated cone with the coordinates reordered.
A strict objectwise quasi-isomorphism of spans induces
one on cones by the cone cohomology sequence.
If a leg was already surjective, then \(f-g\)
is surjective and its kernel is the ordinary pullback.
The short exact sequence with this kernel identifies
its inclusion into the shifted cone as a quasi-isomorphism.
These statements show that fibrational replacement
of one leg followed by ordinary pullback gives a
homotopy-invariant model and that this agrees with
ordinary pullback on spans having that leg fibrant.
\end{proof}

\begin{theorem}[Sequential homotopy limits]\label{res:Q17}
For a tower \(X_0\leftarrow X_1\leftarrow\cdots\) of
complexes, let \(s:\prod_n X_n\to\prod_n X_n\) carry the
\((n+1)\)-st coordinate to the \(n\)-th by the transition map.
A model for \(\operatorname{holim}_n X_n\) is
\[
 \operatorname{Cone}\bigl(1-s:\prod_n X_n\to\prod_n X_n\bigr)[-1].
\]
If all transition maps are degreewise surjective, the
ordinary inverse limit maps quasi-isomorphically to this
model. In each degree \(q\) there is a natural short
exact sequence
\[
 0\longrightarrow\lim\nolimits^1_n H^{q-1}(X_n)
 \longrightarrow H^q(\operatorname{holim}_n X_n)
 \longrightarrow\lim_n H^q(X_n)\longrightarrow0 ,
\]
where for a sequential system of modules
\(\lim^1\) denotes the cokernel of \(1-s\) on their product.
\end{theorem}
\begin{proof}
Products of modules are exact. Taking the cone therefore
gives an objectwise quasi-isomorphism invariant construction.
For a degreewise surjective tower, given a tuple \((b_n)\),
choose \(a_0=0\) and solve
\(s_n(a_{n+1})=a_n-b_n\) recursively in each degree.
Hence \(1-s\) is degreewise surjective and its kernel
is the ordinary inverse limit. The short exact sequence
of complexes proves the asserted comparison with the shifted
cone. Every tower admits a functorial levelwise
quasi-isomorphism to such a tower. Start with
\(Y_0=X_0\). At each subsequent stage factor the composite
\(X_n\to X_{n-1}\to Y_{n-1}\) as a trivial cofibration
\(X_n\to Y_n\) followed by a fibration \(Y_n\to Y_{n-1}\)
using \ref{res:Q12}. Thus taking the ordinary limit
after this replacement is equivalent to the cone model.
Taking cohomology of the cone exact sequence gives the
displayed short exact sequence, since cohomology commutes
with products and the kernel of \(1-s\) on a product
is the ordinary inverse limit.
\end{proof}

\paragraph{Scope.}
The construction above concerns the projective model structure on
module complexes. The unbounded injective resolutions used for
derived Hom are constructed in the derived-functor chapter. The
existence of an injective model structure on sheaf complexes has
different hypotheses and does not enter these proofs.

\chapter{Faithfully flat descent for modules and complexes}
\label{ch:flat-descent}

Let \(A\to B\) be a faithfully flat map of commutative rings. A
quasi-isomorphism can be checked after this base change because flatness
commutes with cohomology and faithfulness detects the zero module. This
chapter will also reconstruct an object from a strict cocycle and recover
an arbitrary unbounded complex from its augmented Amitsur diagram.

\section{The augmented complex and effective module descent}

\begin{definition}[The augmented Amitsur complex]\label{res:F01}
For an \(A\)-module \(M\), put \(C^{-1}(M)=M\) and
\(C^p(M)=B^{\otimes_A(p+1)}\otimes_A M\) for \(p\geq0\). The
map from degree \(-1\) inserts \(1\), and the map from degree \(p\)
is \(\delta=\sum_{i=0}^{p+1}(-1)^id^i\), where \(d^i\) inserts a
unit in position \(i\), counted from zero. This is a complex.
\end{definition}
\begin{proof}
Insertion in two ordered positions \(i<j\) can also be performed by
first inserting in position \(j-1\) and then in position \(i\).
The signs of the two resulting terms of the composite are opposite.
This cancels every term, including the two equal insertions after the
augmentation.
\end{proof}

\begin{theorem}[Exactness after faithfully flat descent]\label{res:F02}
The augmented Amitsur complex of every \(A\)-module is exact.
\end{theorem}
\begin{proof}
After tensoring on the left with \(B\), there is a contraction that
multiplies the new left factor with the first old factor:
\[
 h(b\otimes b_0\otimes\cdots\otimes b_p\otimes m)
   =bb_0\otimes b_1\otimes\cdots\otimes b_p\otimes m .
\]
In degree zero the right hand side is \(bb_0\otimes m\) in the
base-changed augmented term. We have \(hd^0=1\) and
\(hd^i=d^{i-1}h\) for \(i>0\). Therefore
\(h\delta+\delta h=1\), including \(h\delta=1\) on the augmented
term. By faithful flatness, homology of the original complex
must be zero.
\end{proof}

\begin{definition}[Strict module descent data]\label{res:F03}
Write \(B_2=B\otimes_A B\) and \(B_3=B^{\otimes_A3}\). A descent
datum on a \(B\)-module \(N\) is a \(B_2\)-linear isomorphism
\[
 \theta:N\otimes_A B\longrightarrow B\otimes_A N
\]
such that the pullback along multiplication \(B_2\to B\) is the
identity, and \(\theta_{23}\theta_{12}=\theta_{13}\) over \(B_3\).
Subscripts indicate the positions of the old and new \(N\)-factor.
A morphism is a \(B\)-linear map commuting with this isomorphism.
\end{definition}

\begin{theorem}[Effective descent for modules]\label{res:F04}
The functor \(M\mapsto B\otimes_A M\), with its canonical tensor
position isomorphism, is an equivalence onto modules with the
data of Result~\ref{res:F03}. The inverse sends \(N\) to
\[
 M=\{n\in N:\rho(n)=1\otimes n\},\qquad
 \rho(n)=\theta(n\otimes1).
\]
\end{theorem}
\begin{proof}
The map \(\rho\) is \(B\)-linear for the first factor of \(B\otimes_A N\).
The diagonal condition gives \(\mu\rho=1_N\), where \(\mu\) is
multiplication of scalars. If \(\rho(n)=\sum_j b_j\otimes n_j\),
evaluate the cocycle on \(n\otimes1\otimes1\) to get
\[
 (1\otimes\rho)\rho(n)=\sum_j b_j\otimes1\otimes n_j .
\]
Flatness identifies \(B\otimes_A M\) with the submodule of
\(B\otimes_A N\) whose elements \(z\) satisfy
\((1\otimes\rho)(z)=\sum b\otimes1\otimes n\), with
\(z=\sum b\otimes n\). The identity above places \(\rho(N)\)
in this equalizer. Conversely multiplying the first two factors in
the equalizer relation gives \(\rho(\mu z)=z\), using the \(B\)-linearity
of \(\rho\). Thus the natural map \(B\otimes_A M\to N\) is an isomorphism
with inverse \(\rho\) interpreted in the equalizer.
It respects the cocycle because for \(m\in M\) one has
\(\theta(m\otimes1)=1\otimes m\), and \(B_2\)-linearity determines
the isomorphism on all elements. In the canonical case
\(N=B\otimes_A M_0\), the equalizer is \(M_0\) by
Result~\ref{res:F02} in degrees \(-1,0\).
Both constructions are functorial since a map commuting with
\(\theta\) commutes with \(\rho\) and with \(\mu\).
\end{proof}

\begin{corollary}[Descent of finite projective modules]\label{res:F05}
For an \(A\)-module \(M\), if \(B\otimes_A M\) is finite projective
over \(B\), then \(M\) is finite projective over \(A\).
\end{corollary}
\begin{proof}
Express finitely many generators after base change using finitely many
elements of \(M\). The quotient by the module they generate has zero
base change, hence is zero. Choose a finite free surjection to \(M\).
Its kernel becomes finitely generated after base change, since the
surjection then splits and the kernel is a summand of a finite free
module. Apply the same finite-generation argument to this kernel.
Thus \(M\) is finitely presented.
It is flat because for any injection of \(A\)-modules, the kernel of
its tensor with \(M\) becomes zero over \(B\).
For completeness, a finitely presented flat module is projective.
At every localization at a prime, lifts of a residue-field basis
give a free surjection by Nakayama. Its finite kernel has zero
residue since flatness makes tensoring the sequence injective on
the left, so the kernel vanishes by Nakayama. A finite presentation
clears the denominators of this basis and its inverse on a
distinguished neighbourhood of the prime. The chosen finite free
surjection to \(M\) therefore splits over a finite distinguished
cover \(D(f_i)\). By finite presentation, clear denominators in the
local splittings to get maps \(\sigma_i:M\to A^n\) whose composites
with \(A^n\to M\) are \(f_i^{r_i}\mathrm{id}_M\) for sufficiently
large \(r_i\). These powers generate the unit ideal. A linear
combination of the \(\sigma_i\) gives a global splitting.
\end{proof}

\section{The unbounded descent calculation}

\begin{proposition}[Strict descent data for complexes]\label{res:F06}
Chain complexes of \(A\)-modules are equivalent to complexes of
\(B\)-modules equipped with the strict data of Result~\ref{res:F03}
by chain maps. The equivalence also identifies their localizations
at quasi-isomorphisms.
\end{proposition}
\begin{proof}
Apply Result~\ref{res:F04} in each degree. All differentials descend
as morphisms of the descent data. Their composites are zero by the
faithfulness of the equivalence. Tensoring with \(B\) commutes with
cohomology and reflects the zero module, so a map is a
quasi-isomorphism exactly when its base change is one.
The two functors and their comparison isomorphisms consequently
descend to the stated localizations.
\end{proof}

For a cochain complex \(X\) and an integer \(p\geq0\), put
\(C^{p,q}(X)=B^{\otimes_A(p+1)}\otimes_A X^q\).
Horizontal and vertical differentials commute before taking total
signs. Use the differential \(\delta+(-1)^pd_X\) on column \(p\),
and the product over \(p\) in each total degree.
This uses an actual product of modules, not a tensor product with
an already infinite product.

\begin{theorem}[Unbounded Amitsur totalization]\label{res:F07}
For every cochain complex \(X\) the natural augmentation is a
quasi-isomorphism
\[
       X\longrightarrow \operatorname{Tot}^{\Pi} C^{\bullet,\bullet}(X).
\]
The target is also the homotopy limit in \(D(A)\) of totalizations
truncated to the first finitely many columns.
\end{theorem}
\begin{proof}
For bounded-below \(X\), add the augmented column \(p=-1\) with
vertical sign \(-1\). The augmented total complex has acyclic rows
by Result~\ref{res:F02}. The lower bound on the vertical degree
and the finite number of columns in each total degree justify
the row filtration argument of \cite[B01--B04]{ha-core}. The
augmented total complex is the cone of the augmentation, with
column \(p=-1\) equal to \(X[1]\) and carrying differential \(-d_X\).
Its acyclicity proves the first assertion for bounded-below \(X\).

For general \(X\) put \(X_n=\tau_{\geq-n}X\) for \(n\geq0\), using
the quotient by boundaries at degree \(-n\), zero below and the
original modules above. The canonical tower maps \(X_{n+1}\to X_n\)
are degreewise surjective, the inverse limit is \(X\), and these
maps eventually become the identity in each fixed degree.
The augmentation maps for \(X_n\) form a map of towers.
For each column and each internal degree, tensoring with the indicated
finite tensor power of \(B\) commutes with this inverse limit because
the tower in that fixed degree is eventually constant.
Products commute with inverse limits, so the limit of the product
totalizations of the \(X_n\) is exactly the displayed product
totalization. Their transitions are degreewise surjective because
they are products of surjections.

For a tower \(Y_n\) with degreewise surjective transitions \(u_n\),
there is a short exact sequence of complexes
\[
 0\to\varprojlim_n Y_n\to\prod_{n\geq0}Y_n
 \xrightarrow{1-\mathrm{shift}}\prod_{n\geq0}Y_n\to0,
 \qquad (1-\mathrm{shift})(y)_n=y_n-u_n(y_{n+1}).
\]
Surjectivity follows by solving from \(n=0\) onwards. Products of
module complexes preserve quasi-isomorphisms because products of
surjections are surjective and kernels are products of kernels.
Apply this short exact sequence to the map of towers
\(X_n\to\operatorname{Tot}^{\Pi} C(X_n)\). The two product maps are
quasi-isomorphisms, so the long exact sequence gives the claim on
their inverse limits. This argument does not require \(B\otimes_A-\)
to preserve infinite products.

Finally truncate at horizontal column \(p=n\), by quotienting the
subcomplex of later columns. The projections between truncated
totalizations are surjective termwise, and the ordinary inverse
limit is the product totalization. The same short exact sequence
identifies that limit with the shifted cone of \(1-\mathrm{shift}\),
the sequential homotopy-limit model of Result~\ref{res:Q17}.
\end{proof}

\begin{corollary}[Descent of perfect-coefficient mapping complexes]\label{res:F08}
Let \(P\) be a bounded complex of finite projective \(A\)-modules
and \(X\) any complex. With \(B_p=B^{\otimes_A(p+1)}\) there is an
isomorphism in \(D(A)\)
\[
 \mathrm{RHom}_A(P,X)\simeq
 \operatorname{Tot}^{\Pi}_{p\geq0}
 \mathrm{RHom}_{B_p}(B_p\otimes_A P,B_p\otimes_A X),
\]
where the right hand side uses the cofaces of Result~\ref{res:F01}
on the compatible ordinary Hom models.
\end{corollary}
\begin{proof}
Bounded complexes of projectives are K-projective by the finite
filtration by degree tails. Thus the ordinary Hom complexes
compute each term. A finite projective module is a direct summand
of a finite free one. This gives
\(B_p\otimes_A\operatorname{Hom}_A(P^j,X^q)
=\operatorname{Hom}_{B_p}(B_p\otimes P^j,B_p\otimes X^q)\)
canonically, with equality here denoting the base-change isomorphism.
Each Hom degree uses only finitely many \(j\), so this remains
compatible with the Hom differential, multiplication and unit
on the base rings. Apply Result~\ref{res:F07} to
\(\operatorname{Hom}_A^\bullet(P,X)\).
\end{proof}

\begin{example}[Localization covers and the need for faithfulness]\label{res:F09}
If \(f_1,\ldots,f_s\) generate the unit ideal of \(A\), then
\(B=\prod_i A_{f_i}\) is faithfully flat and its Amitsur diagram
is the full ordered affine-intersection diagram of the cover by
\(D(f_i)\). In contrast, the flat map \(\mathbb Z\to\mathbb Q\)
does not detect the nonzero complex \(\mathbb Z/\ell\) in degree
zero for a prime \(\ell\).
\end{example}
\begin{proof}
The product is finite, hence a direct sum of flat modules as an
\(A\)-module. If \(m\) vanishes at each localization then it is
killed by a positive power of each \(f_i\). Those powers generate
the unit ideal, so \(m=0\). In particular base change reflects
zero modules, proving faithful flatness. The tensor product
distributes over the finite ring factors, and the tensor product
of the corresponding localizations is localization at the
product \(f_{i_0}\cdots f_{i_p}\). In the component indexed by a
longer ordered tuple, insertion of \(1\) restricts from the
intersection obtained by omitting that entry to the full intersection.
Finally \(\ell\) acts invertibly over \(\mathbb Q\), so tensoring
\(\mathbb Z/\ell\) with \(\mathbb Q\) gives zero.
\end{proof}

These are strict and explicit forms of descent. The general
homotopy-coherent formulation replaces ordinary compatibility by
a descent datum valued in a homotopy limit of higher categories.
Derived stacks and hypercovers use that framework
\cite[\S1.2.12]{hag}. The product calculation above already records
the higher Amitsur degrees and their total signs for affine
complexes.

\chapter{Derived affine intersections and cotangent complexes}
\label{ch:derived-affines}

Let \(k\) be a field of characteristic zero. In this chapter algebras
are commutative dg \(k\)-algebras in cochain degrees \(\leq0\),
with \(ab=(-1)^{|a||b|}ba\). An ordinary \(k\)-algebra is placed
in degree zero. All constructions use graded commutativity and
specified dg algebra maps. The affine convention reverses algebra
maps. We call the object represented by \(B\) a derived affine and
regard quasi-isomorphic algebras as equivalent. We prove the affine
calculations underlying this notation, following the scope of
\cite[\S\S2.2 and 4.1]{toen}.

\section{Semifree resolutions of commutative dg algebras}

\begin{definition}[Semifree algebra extension]\label{res:A01}
A semifree extension of \(A\) is a graded-commutative polynomial
extension on homogeneous variables adjoined in stages, with
the differential of a variable in the algebra at earlier stages.
The free graded-commutative algebra is polynomial in even variables
and exterior in odd variables. If prescribed differentials are
cycles of the required degrees, the unique graded derivation extending
them is a differential.
\end{definition}
\begin{proof}
The Koszul sign rule gives a well-defined derivation on the free
graded-commutative algebra. In particular for an odd variable \(x\),
\(d(x^2)=d(x)x-xd(x)=0\) since \(d(x)\) is even.
The cross terms in applying a degree-one derivation twice cancel.
Thus \(d^2\) vanishes on every monomial if it vanishes on the base
and each variable. Each specified cycle ensures that condition.
\end{proof}

\begin{theorem}[Semifree resolution and lifting]\label{res:A02}
Every map \(A\to B\) admits a factorization
\(A\to P\twoheadrightarrow B\) with \(P\) semifree over \(A\)
and the second arrow a degreewise surjective quasi-isomorphism.
A map from a semifree extension lifts through any such surjective
quasi-isomorphism of \(A\)-algebras.
\end{theorem}
\begin{proof}
In degree zero adjoin cycles mapping to all elements of \(B^0\).
For each \(b\in B^n\), \(n<0\), first adjoin a cycle \(v_b\)
in degree \(n+1\) mapping to \(db\), then \(u_b\) in degree \(n\)
with \(du_b=v_b\), mapping to \(b\). This makes the map onto \(B\)
surjective in every degree. At the next stage, for every homogeneous
cycle \(z\) in the kernel adjoin a generator of degree \(|z|-1\)
with differential \(z\) and image zero. Repeat countably many times.
Every element is a finite polynomial using finitely many stages.
Thus every kernel cycle in the union occurred at a stage and became
a boundary at the next. The kernel is acyclic and the surjection is
a quasi-isomorphism.

For lifting through \(C\twoheadrightarrow D\), choose the lift of a
generator \(x\) after all generators in \(dx\) have been lifted.
A candidate lift \(c\) has \(dc-\widetilde{dx}\) a cycle in the kernel.
Subtract a primitive of that cycle to get the required lift. Proceed
through the stages. In degree zero the candidate is already a cycle
because the target is nonpositive. The construction respects the base
and is defined on the union.
\end{proof}

\begin{lemma}[K-flatness of a semifree algebra]\label{res:A03}
A semifree \(P\) over \(A\), regarded as a dg \(A\)-module,
is K-flat.
\end{lemma}
\begin{proof}
At one adjoining stage use the increasing filtration by total degree
in the new variables. Differentiating such a variable removes it and
gives a coefficient from the previous stage. Each graded quotient of
the filtration is therefore a direct sum of shifts of the previous
stage as a dg module. The inclusions are split as graded modules by
the monomial bases. Tensor with an acyclic dg \(A\)-module. Starting
from \(A\), the quotients after tensoring are acyclic and the graded
splittings preserve exactness of successive sequences after tensoring.
The long exact sequence and then an exact filtered union prove
acyclicity at that adjoining stage. Induct through the stages, using
filtered exactness at their unions.
\end{proof}

\section{Derived base change and intersection}

\begin{proposition}[Tensor model for a derived affine base change]\label{res:A04}
For \(A\to B\) and \(A\to C\), a model of the derived tensor product
is \(P\otimes_A C\) for a semifree surjection \(P\to B\) as in
Result~\ref{res:A02}. A change of models gives a zigzag of algebra
quasi-isomorphisms. Its zeroth cohomology is
\[
 H^0(P\otimes_A C)
   \cong H^0(B)\otimes_{H^0(A)} H^0(C).
\]
\end{proposition}
\begin{proof}
Choose semifree surjective resolutions of both \(B\) and \(C\).
Tensoring the two resolutions computes the derived tensor product of
dg \(A\)-modules by Result~\ref{res:A03}. Replacing the second one
by \(C\) is then a quasi-isomorphism by K-flatness of \(P\).
For two choices for \(B\), lift the augmentation from one to the
other by Result~\ref{res:A02}; the lifted map is a quasi-isomorphism
by the comparison to \(B\). Tensor with the semifree model for \(C\)
and use its K-flatness. This gives the stated zigzag of maps of
commutative dg algebras, and the argument is symmetric in the two
inputs. We use this construction as the derived affine fibre product
model.

Since all degrees are nonpositive, degree zero of the tensor is
generated by \(P^0\otimes_{A^0}C^0\), and its boundaries are exactly
the images obtained by using \(dP^{-1}\) in the first factor or
\(dC^{-1}\) in the second. The balancing relations in degree zero
come from \(A^0\). Its boundaries act as zero after the two
quotients, so the result factors over \(H^0(A)\), giving the formula.
\end{proof}

\begin{example}[The Koszul model of a derived zero locus]\label{res:A05}
Let \(A\) be ordinary and let \(f_1,\ldots,f_r\in A\).
Regard \(A\) as an algebra over \(k[t_1,\ldots,t_r]\) by \(t_i\mapsto f_i\).
The derived fibre over the zero section has algebra
\[
     A[e_1,\ldots,e_r],\qquad |e_i|=-1,\qquad de_i=f_i .
\]
If \(f_i\) form an \(A\)-regular sequence then this algebra resolves
\(A/(f_1,\ldots,f_r)\). In that case the derived self-intersection
of the closed locus over \(A\) has exterior generators in degree
\(-1\) and zero differential over \(A/(f_1,\ldots,f_r)\).
\end{example}
\begin{proof}
The analogous algebra with differential \(de_i=t_i\) is the Koszul
resolution of \(k\) over the polynomial ring, because the variables
are regular. The exactness is the Koszul resolution calculation
\cite[RK-07]{ca-core}. It is semifree, so tensoring with \(A\)
gives the derived fibre by Result~\ref{res:A04}.
For a regular sequence in \(A\), the same Koszul exactness applies
over \(A\). The displayed algebra is then a semifree resolution of
the quotient over \(A\). Tensor with the quotient, which sends
each \(f_i\) to zero. For example the derived self-intersection of
the origin in \(\mathbb A^1_k\) has \(H^0=k\) and \(H^{-1}=k e_1\).
\end{proof}

\section{The cotangent calculation}

Use \(\partial\) for a K\"ahler differential and \(d\) for the
cochain differential. For a dg \(P\)-module, the usual graded Hom
complex is formed using the Koszul rule for the degree of a
homogeneous linear map.

\begin{lemma}[Differentials on a semifree algebra]\label{res:A06}
For \(A\to P\) semifree, the dg module \(\Omega_{P/A}\) is free
as a graded \(P\)-module on \(\partial x\) for the variables \(x\),
where \(|\partial x|=|x|\) and \(d(\partial x)=\partial(dx)\).
It is a semifree \(P\)-module. Its base change to any dg \(P\)-algebra
is again semifree. The complex
\(\operatorname{Hom}^{\bullet}_P(\Omega_{P/A},M)\) identifies with
graded \(A\)-derivations \(P\to M\) with differential the graded
commutator with \(d\).
\end{lemma}
\begin{proof}
The graded Leibniz relation extends arbitrary values on the free
generators uniquely to a graded derivation, and it respects
the symmetric and exterior relations. Hence the graded module of
differentials has the stated basis. The relation requiring
\(\partial\) to commute with the cochain differential determines
the formula for \(d\) on that basis. The derivative of \(dx\) uses
only earlier generators, so the corresponding filtration gives
semifreeness and is preserved by base change. The Hom--derivation
identification now commutes with its differential by evaluating
the commutators on the same basis.
\end{proof}

\begin{lemma}[A relative polynomial path comparison]\label{res:A07}
Let \(Q\twoheadrightarrow B\) be a surjective quasi-isomorphism of
nonpositive dg algebras. Let \(A\to P\) be semifree with a map
\(P\to B\).
Any two \(A\)-algebra maps \(P\to Q\) over \(B\) are joined by a
polynomial dg path over \(B\). Consequently their induced maps
\[
 B\otimes_P\Omega_{P/A}\longrightarrow B\otimes_Q\Omega_{Q/A}
\]
are chain homotopic as maps of dg \(B\)-modules.
\end{lemma}
\begin{proof}
Let \(k[t,dt]\) have \(|t|=0\), \(|dt|=1\) and differential
\(d(t)=dt\), with \(dt\) exterior. Endpoints set \(t=0,1\) and
\(dt=0\). The inclusion of constants and either evaluation are
mutual homotopy inverses as complexes by polynomial integration.
On a homogeneous coefficient \(y\), the integration of
\(y\,g(t)dt\) has sign \((-1)^{|y|}\).
The identity \(dH+Hd=\operatorname{ev}_1-\operatorname{ev}_0\)
for this integration follows by the fundamental formula for a
polynomial derivative and the opposite signs on the two terms
containing \(d y\). Integers are invertible in \(k\).

Take the relative path algebra
\[
 T=\tau_{\leq0}\bigl(Q[t,dt]\times_{B[t,dt]} B\bigr),
\]
where the truncation keeps the kernel of \(d\) in degree zero and
all degrees below zero. This is a subalgebra. Write \(K\) for the
acyclic kernel of \(Q\to B\). Before truncation, evaluation at zero
has section given by constants and kernel
\(K\otimes_k\ker(k[t,dt]\to k)\) for evaluation at zero.
This kernel is contractible using integration from zero to \(t\).
The good truncation preserves the cohomology comparison to \(Q\).
The projection \(Q\times_B Q\to Q\) is also a quasi-isomorphism
since its kernel is \(K\). Therefore
\(\operatorname{ev}_{01}:T\to Q\times_B Q\) is a quasi-isomorphism.
It is surjective in negative degrees using linear interpolation
of endpoints. In degree zero, if endpoints differ by \(w\in K^0\),
pick \(z\in K^{-1}\) with \(dz=w\). Then the path
\(q_0+tw-z\,dt\) is closed and has the required endpoints.
Use the lifting part of Result~\ref{res:A02} for \(P\) to get the
path.

To compare on differentials, differentiate in the \(Q\)-variable
of the path, leaving \(t,dt\) constant. This maps differentials
into \(\Omega_{Q/A}[t,dt]\). After extension along \(Q\to B\), the
original path has constant image in \(B[t,dt]\). The derivative
therefore induces a dg \(B\)-linear map into
\((B\otimes_Q\Omega_{Q/A})[t,dt]\). Its endpoint maps are the two
maps in the statement. Integration gives a graded \(B\)-linear
chain homotopy between them.
\end{proof}

\begin{theorem}[The affine cotangent complex]\label{res:A08}
For a semifree surjective resolution \(P\to B\) over \(A\), the
object
\[
       L_{B/A}=B\otimes_P\Omega_{P/A}
\]
is independent of the resolution up to canonical isomorphism in
the homotopy category of dg \(B\)-modules. It is represented by
a K-projective, K-flat dg \(B\)-module.
\end{theorem}
\begin{proof}
For two resolutions lift the maps over \(B\) in both directions
by Result~\ref{res:A02}. Their composites and the identities
agree after augmentation. Result~\ref{res:A07} proves that their
differential maps are mutual homotopy inverses and independent
of the chosen lifts. Composing three comparisons respects these
homotopy classes, giving the canonical isomorphisms.
The dg module in the statement is semifree by
Result~\ref{res:A06}. A semifree dg module is K-flat by its
filtration with shifted free quotients. To see K-projectivity,
a cocycle map from it into an acyclic dg module can be made
nullhomotopic successively on generators. On a new generator,
subtract the already specified homotopy on its differential.
The resulting prescribed boundary equation has a cycle on
its left, so acyclicity supplies a primitive. The Koszul sign
determines the module-linear extension of the homotopy.
Apply this in every degree by shifting the target. Thus the
Hom complex into an acyclic dg module is acyclic.
\end{proof}

\begin{theorem}[Transitivity and base change of the cotangent complex]\label{res:A09}
For \(A\to B\to C\) as above there is a natural triangle in \(D(C)\)
\[
 C\otimes_B^{\mathbf L}L_{B/A}\longrightarrow L_{C/A}
 \longrightarrow L_{C/B}\longrightarrow
 (C\otimes_B^{\mathbf L}L_{B/A})[1].
\]
For \(A\to A'\), if \(P\to B\) is semifree and
\(B'=A'\otimes_A P\) is the derived base-change model, then
\[
       L_{B'/A'}\simeq B'\otimes_P \Omega_{P/A}.
\]
In the flat ordinary base-change case, the second formula identifies
with extension of scalars on \(L_{B/A}\) along \(B\to A'\otimes_A B\).
\end{theorem}
\begin{proof}
Choose \(P\to B\) semifree over \(A\), then resolve \(C\) by \(Q\)
semifree over \(P\) along \(P\to B\to C\), using
Result~\ref{res:A02}. Concatenation makes \(Q\) semifree over \(A\).
The variable bases give a graded-split exact sequence of dg modules
\[
 0\longrightarrow Q\otimes_P\Omega_{P/A}
 \longrightarrow \Omega_{Q/A}
 \longrightarrow \Omega_{Q/P}\longrightarrow0 .
\]
The first submodule is a subcomplex since the differential of
a \(P\)-variable uses only \(P\)-variables. The quotient is specified
by setting \(\partial x=0\) for each \(P\)-variable \(x\).
Tensor with \(C\) over \(Q\); graded splitness preserves exactness.
The first resulting term equals
\(C\otimes_B(B\otimes_P\Omega_{P/A})\), with the factor in parentheses
K-flat by Result~\ref{res:A08}. The middle term computes \(L_{C/A}\).
For the last term \(B\otimes_P Q\) is semifree over \(B\) and
quasi-isomorphic to \(C\): the map from \(Q\) is a quasi-isomorphism
by Result~\ref{res:A03}, and \(Q\to C\) is a quasi-isomorphism.
The map \(B\otimes_PQ\to C\) is surjective since \(Q\to C\) is.
Relative differentials commute with extension of scalars by the
variable formula. Thus the last term computes \(L_{C/B}\).
The exact sequence supplies the triangle. To compare two resolved
diagrams, lift first on \(P\) and then on the variables of \(Q\),
over the prescribed maps on \(B,C\). Such a comparison commutes
with the injection and quotient maps on differentials. For two
choices of comparison to \(P'\to Q'\), choose a relative path
between the two maps on \(P\) over \(B'\) by
Result~\ref{res:A07}. The map \(P'\to Q'\) carries it to the
relative path algebra of \(Q'\to C'\), which now has a \(P\)-algebra
structure. The endpoint map to \(Q'\times_{C'}Q'\) is a surjective
quasi-isomorphism. Lift through it using the semifree \(Q\)-cells
over \(P\). This yields a path of diagram comparisons with the
chosen path on \(P\). Differentiation and integration as in
Result~\ref{res:A07} give compatible homotopies on the short exact
sequences, proving independence and naturality in the derived
category (with extension of scalars on the target modules).

Finally \(A'\otimes_A P\) is itself semifree over \(A'\);
the relative variables give the claimed expression on the same
generators and differential. For flat ordinary base change the
map \(A'\otimes_A P\to A'\otimes_A B\) is a surjective
quasi-isomorphism. Using this resolution and the associativity
of tensor gives the final identification.
\end{proof}

\begin{proposition}[Square-zero tangent calculation]\label{res:A10}
Let \(M\) be a nonpositive dg \(B\)-module. Augmented maps
\(P\to B\ltimes M\) of \(A\)-algebras over the fixed augmentation
\(P\to B\) are naturally the degree-zero cycles in
\(\operatorname{Hom}_B^\bullet(L_{B/A},M)\) for the model of
Result~\ref{res:A08}. Their classes modulo polynomial dg homotopy
over \(B\) are \(H^0\mathrm{RHom}_B(L_{B/A},M)\).
All groups \(H^i\mathrm{RHom}_B(L_{B/A},M)\) are computed by the
same Hom complex.
\end{proposition}
\begin{proof}
The multiplication on \(B\ltimes M\) has square-zero ideal \(M\).
The difference of an augmented map from the map with zero
\(M\)-component is therefore an \(A\)-derivation by the Leibniz rule.
Commutation with the differential is its degree-zero cocycle
condition. Result~\ref{res:A06} gives the correspondence, and
K-projectivity in Result~\ref{res:A08} computes the derived Hom.

For the assertion on classes, a polynomial path is a closed
degree-zero derivation with values in \(M[t,dt]\) over the
constant augmentation, with a good truncation if necessary.
Integration as in Result~\ref{res:A07} makes the endpoints differ
by a Hom boundary. Conversely, if two cocycles \(f_0,f_1\) have
\(f_1-f_0=d_{\rm Hom}h\) for \(h\) of degree \(-1\), the path
\(f_0+t(f_1-f_0)-h\,dt\), using the tensor sign convention, is
closed of degree zero. Here the expression means the map into
\(M\otimes_k k[t,dt]\) obtained by tensoring graded maps with the
specified forms. It is well-defined by the graded tensor sign
on homogeneous arguments. Its degree-zero values on closed
elements are closed, so it factors through the good truncation.
It therefore gives the required dg algebra path. This proves
the quotient identification.
\end{proof}

\section{Complete intersections and compatibility with descent}

\begin{example}[The conormal complex and its tangent space]\label{res:A11}
Let \(B=S/(f_1,\ldots,f_r)\) with \(S=k[x_1,\ldots,x_s]\) and
\(f_i\) regular in \(S\). Then \(L_{B/k}\) is represented by
\[
    [\,B^r\xrightarrow{(\partial f_i/\partial x_j)}B^s\,]
\]
in degrees \(-1,0\), with column \(i\) the derivative of \(f_i\).
In particular \(H^1\mathrm{RHom}_B(L_{B/k},B)\) is the first-order
deformation space of the marked algebra in
\cite[T04]{ca-topics}.
\end{example}
\begin{proof}
Use the semifree Koszul algebra over \(S\) of Result~\ref{res:A05}
and regard the variables \(x_j\) as its first cells over \(k\).
The differential on \(\partial e_i\) is \(\partial f_i\);
the differentials on \(\partial x_j\) vanish. Tensor with \(B\)
to get precisely the two-term complex.
Taking Hom into \(B\) gives in degree one the cokernel of
the transposed Jacobian. The explicit square-zero presentation
calculation in \cite[T02--T04]{ca-topics} identifies that cokernel
with the deformation classes over the dual numbers.
\end{proof}

\begin{proposition}[Faithfully flat locality of the affine tangent calculation]\label{res:A12}
Let \(A\) be an ordinary \(k\)-algebra, \(A\to C\) faithfully
flat, and \(B\) an ordinary \(A\)-algebra. Assume the affine
cotangent complex \(L_{B/A}\) is represented by a bounded complex
\(P_L\) of finite projective \(B\)-modules. Set
\(C_p=C^{\otimes_A(p+1)}\) and \(B_p=B\otimes_A C_p\).
For any complex \(M\) of \(B\)-modules there is a comparison
\[
 \mathrm{RHom}_B(L_{B/A},M)
 \simeq \operatorname{Tot}^{\Pi}_{p\geq0}
     \mathrm{RHom}_{B_p}(L_{B_p/C_p}, B_p\otimes_B M).
\]
When all local tangent calculations are made in the ordinary
degree-zero case, this compares their classes and higher
cohomology by the indicated full total complex.
\end{proposition}
\begin{proof}
Base change \(B\to B\otimes_A C\) is faithfully flat.
Each \(C_p\) is flat over \(A\). The base-change part of
Result~\ref{res:A09} uses a single semifree model over \(A\)
and its compatible scalar extensions to the whole Amitsur diagram.
Thus \(B_p\otimes_B P_L\) represents \(L_{B_p/C_p}\).
A comparison from \(P_L\) to the cofibrant cotangent model can
be chosen as a chain map since \(P_L\) is K-projective, and its
flat scalar extensions commute with the cofaces. Apply
Result~\ref{res:F08} over \(B\) to \(P_L\) and \(M\), interpreting
the Hom complexes using these common models.
\end{proof}

The affine calculations here use a characteristic-zero commutative
dg model. In general, derived stacks involve functors of points,
homotopy-coherent descent and representability. These broader
constructions, as well as the all-characteristic simplicial
setting, are developed in \cite[\S\S1.2.12 and 2.2.3]{hag} and
\cite[\S\S2 and 4]{toen}.

\backmatter
\begingroup
\small
\begin{thebibliography}{20}
\setlength{\itemsep}{0.2\baselineskip}
\setlength{\parsep}{0pt}
\bibitem{p103} MathGenome. \emph{Algebra Core} (P1-03). v0.22.0,
  Chapters 2, 4, 6--7, 10--11.
\bibitem{ha-core} MathGenome. \emph{Advanced Homological Algebra}.
  Independently audited notes, 19 November 2026 snapshot.
\bibitem{ca-core} MathGenome. \emph{Advanced Commutative Algebra:
  Differentials, Smoothness, and Local Algebra}. Independently audited
  notes, 19 November 2026 snapshot.
\bibitem{ca-topics} MathGenome. \emph{Topics in Commutative Algebra}.
  20 November 2026, results L01--L11, D01--D09 and T01--T04.
\bibitem{keller} B. Keller. On differential graded categories.
  In \emph{Proceedings of the International Congress of Mathematicians},
  Madrid, Vol. II, 2006, pp. 151--190. \url{https://arxiv.org/abs/math/0601185}.
\bibitem{yekutieli} A. Yekutieli. \emph{Derived Categories}.
  Cambridge University Press, 2020, Chapters 3, 12 and 13.
\bibitem{hovey} M. Hovey. \emph{Model Categories}.
  Mathematical Surveys and Monographs 63, American Mathematical Society,
  1999. Author manuscript archived with the accompanying corrections.
\bibitem{hovey-errata} M. Hovey. Corrections to \emph{Model Categories}.
  Author's errata.
\bibitem{riehl} E. Riehl. \emph{Categorical Homotopy Theory}.
  New Mathematical Monographs 24, Cambridge University Press, 2014.
\bibitem{st-derived} The Stacks Project Authors. \emph{Derived Categories}.
  \url{https://stacks.math.columbia.edu/tag/05QI}.
\bibitem{st-sheaf} The Stacks Project Authors. \emph{Cohomology of Sheaves}.
  \url{https://stacks.math.columbia.edu/tag/01DW}.
\bibitem{st-duality} The Stacks Project Authors. \emph{Duality for Schemes}.
  \url{https://stacks.math.columbia.edu/tag/0DWE}.
\bibitem{st-descent} The Stacks Project Authors. \emph{Descent for modules}
  and \emph{Descent of quasi-coherent sheaves}, Chapter 35, \S\S3--5.
  \url{https://stacks.math.columbia.edu/tag/023F}.
\bibitem{toen} B. To\"en. Derived algebraic geometry.
  \emph{EMS Surveys in Mathematical Sciences} 1 (2014), 153--240.
  \url{https://arxiv.org/abs/1401.1044}.
\bibitem{hag} B. To\"en and G. Vezzosi.
  \emph{Homotopical Algebraic Geometry II}.
  Memoirs of the American Mathematical Society 193 (2008), no.\ 902.
  \url{https://arxiv.org/abs/math/0404373}.
\end{thebibliography}
\endgroup

\end{document}
